% ============================================================
%  tilvista_doku.tex  --  Hauptdokument
%
%  Kompilierung (3 Durchlaeufe fuer TOC + Index):
%    pdflatex tilvista_doku.tex
%    makeindex tilvista_doku.idx
%    pdflatex tilvista_doku.tex
%    pdflatex tilvista_doku.tex
%
%  Verzeichnisstruktur:
%    tilvista_doku.tex        <- dieses Dokument
%    graphics/
%      icon.png               <- App-Icon (PNG, mind. 64x64)
%    src/
%      packages/
%        packages.tex         <- alle \usepackage-Befehle
%        macros.tex           <- Farben, Stile, Makros, Seitenstil
%      chapter1_allgemein.tex
%      chapter2_aleavue.tex
%      chapter3_sattumapic.tex
%      chapter4_madoludus.tex
%      chapter5_workers.tex
%      chapter6_zusammenfassung.tex
% ============================================================
\documentclass[
  12pt,
  a4paper,
  twoside,
  headings=normal,
  parskip=half
]{scrreprt}

% -- Pakete und Makros einbinden -----------------------------
% Reihenfolge ist wichtig: packages.tex laedt alle Pakete,
% macros.tex definiert Farben/Stile/Befehle die diese Pakete
% voraussetzen. Beide vor \begin{document}.
% ============================================================
%  packages.tex  --  alle eingebundenen LaTeX-Pakete
%
%  WARUM KEIN \usepackage{babel} MIT NGERMAN?
%  ------------------------------------------
%  Babel ist das Standard-Lokalisierungspaket fuer LaTeX und
%  liefert u.a. deutsche Trennmuster, Anpassung von
%  Kapitelbeschriftungen ("Kapitel" statt "Chapter") und
%  korrekte Silbentrennung.
%
%  Auf diesem System (TeX Live 2023/Debian) fehlt das Paket
%  "babel-german" (normalerweise als separates CTAN-Paket
%  installiert). Ohne es kennt babel die Option "ngerman"
%  nicht und bricht den Build ab.
%
%  Alternativen und warum sie hier nicht greifen:
%   - \usepackage{ngermanb}   --> erwartet ngermanb.ldf, fehlt
%   - \usepackage[german]{babel} --> braucht german.ldf, fehlt
%   - polyglossia  --> nur fuer XeLaTeX/LuaLaTeX, wir nutzen pdfLaTeX
%
%  Was wir stattdessen tun:
%   - inputenc + fontenc:  korrekte UTF-8-Eingabe und T1-Ausgabe
%     (Umlaute in Fliesstextbloecken funktionieren damit einwandfrei)
%   - csquotes:  deutsche Anführungszeichen per \enquote{}
%   - Listings: extendedchars=false (verhindert UTF-8-Fehler
%     INNERHALB von lstlisting-Umgebungen; Kommentare dort
%     deshalb in ASCII-Umschreibung)
%   - Silbentrennung laeuft ohne babel auf Englisch; das ist
%     bei einem technischen Dokument mit vielen Code-Bezeichnern
%     kaum sichtbar.
%
%  Sobald babel-german auf dem System verfuegbar ist, reicht
%  es, die beiden auskommentierten Zeilen unten zu aktivieren
%  und die "\usepackage{csquotes}"-Zeile zu behalten.
%
%  Zum Aktivieren (wenn babel-german installiert):
%    \usepackage[ngerman]{babel}
%    \usepackage{csquotes}   % arbeitet dann mit babel zusammen
% ============================================================

% -- Zeichenkodierung ----------------------------------------
% inputenc: Eingabedatei wird als UTF-8 gelesen.
%   Ohne dieses Paket muss jeder Umlaut als Escape-Sequenz
%   geschrieben werden (\"a statt ae, etc.).
\usepackage[utf8]{inputenc}

% fontenc: Ausgabe-Zeichensatz T1 (8-bit).
%   T1 enthaelt alle westeuropaeischen Zeichen als echte
%   Glyphen -- Silbentrennung an Umlauten funktioniert damit
%   korrekt. Ohne T1 wuerde z.B. "Oe" als zwei Zeichen
%   behandelt und nicht sauber getrennt.
\usepackage[T1]{fontenc}

% csquotes: typografisch korrekte Anführungszeichen.
%   \enquote{Text} erzeugt "Text" (je nach Locale).
%   Ohne babel-Kopplung nutzen wir den Stil manuell:
%   \MakeOuterQuote{"} wäre eine Alternative, aber
%   \enquote ist sicherer und semantisch klarer.
\usepackage{csquotes}

% -- Schriften -----------------------------------------------
% palatino: Serifenschrift (Fliesstext).
%   Besser lesbar als Computer Modern auf Bildschirm und
%   Ausdruck; wird von Qt-Dokumentationen aehnlich verwendet.
\usepackage{palatino}

% courier: Monospace-Schrift fuer Code-Listings.
%   Ist auf praktisch jedem TeX-System vorhanden und
%   unterstuetzt den T1-Zeichensatz vollstaendig.
\usepackage{courier}

% helvet: serifenlose Schrift (Helvetica) fuer Ueberschriften,
%   skaliert auf 92 % der normalen Groesse (wirkt proportionaler).
\usepackage[scaled=0.92]{helvet}

% -- Seitenlayout --------------------------------------------
% geometry: Seitenraender fein einstellen.
%   inner/outer statt left/right fuer zweiseitigen Druck:
%   "inner" ist der Bundsteg (nahe am Heftrand), "outer" der
%   Aussensteg. headheight=15pt verhindert eine fancyhdr-Warnung.
\usepackage[
  top=2.8cm, bottom=3.2cm,
  inner=3.0cm, outer=2.4cm,
  headheight=15pt
]{geometry}

% fancyhdr: individuelle Kopf- und Fusszeilen.
%   Standardmaessig zeigt LaTeX nur Seitennummern;
%   fancyhdr erlaubt Kapiteltitel + Icon in der Kopfzeile.
\usepackage{fancyhdr}

% -- Farben --------------------------------------------------
% xcolor mit dvipsnames: vordefinierte Farbnamen (z.B. BrickRed)
%   und table-Option fuer gefaerbte Tabellenzeilen.
\usepackage[dvipsnames,table]{xcolor}

% -- Schriften (Ergaenzung) ----------------------------------
% microtype: verbessert den Textumbruch durch optischen
%   Randausgleich und Zeichenabstandsanpassung (nur pdfLaTeX).
%   Rein typografisch, kein sichtbarer Effekt auf Struktur.
\usepackage{microtype}

% -- Grafik --------------------------------------------------
% graphicx: \includegraphics fuer PNG/PDF/JPG.
%   Benoetigt fuer das Icon in Titelseite und Kopfzeile.
\usepackage{graphicx}

% -- Code-Listings -------------------------------------------
% listings: typsetzt Quellcode mit Syntax-Highlighting.
%   Eigene Stile werden in macros.tex definiert.
%   Hinweis: extendedchars=false ist in den Stilen gesetzt,
%   weil listings intern keine echte UTF-8-Unterstuetzung hat
%   und bei Umlauten in Code-Bloecken sonst abbricht.
\usepackage{listings}

% -- Boxen (Hinweise, Warnungen, Qt-Hintergrund) -------------
% tcolorbox mit skins und breakable:
%   "skins" erlaubt erweiterte Rahmengestaltung,
%   "breakable" erlaubt Boxen ueber Seitenumbrueche hinweg.
\usepackage{tcolorbox}
\tcbuselibrary{skins,breakable}

% -- Tabellen ------------------------------------------------
% booktabs: professionelle horizontale Linien (\toprule etc.)
%   ohne vertikale Linien -- entspricht typografischen Regeln.
\usepackage{booktabs}

% longtable: Tabellen ueber mehrere Seiten.
%   Benoetigt fuer die Kuerzel- und Klassen-Referenz im Anhang.
\usepackage{longtable}

% array: erweiterte Spalten-Definitionen in Tabellen
%   (z.B. feste Spaltenbreite mit p{3cm}).
\usepackage{array}

% -- Listen --------------------------------------------------
% enumitem: feinere Kontrolle ueber Listenabstaende.
%   noitemsep reduziert den Abstand zwischen Listenpunkten
%   fuer kompaktere Darstellung.
\usepackage{enumitem}

% -- Absatzgestaltung ----------------------------------------
% parskip: ersetzt Einrueckung durch halben Absatzabstand.
%   Passt besser zu einem technischen Dokument als klassische
%   Einrueckung; kombiniert mit scrreprt's parskip=half.
\usepackage{parskip}

% -- Index ---------------------------------------------------
% makeidx: erzeugt einen Stichwortindex am Dokumentende.
%   \index{Begriff} im Text, \printindex am Ende.
\usepackage{makeidx}
\makeindex

% -- Hyperlinks und PDF-Metadaten ----------------------------
% hyperref: macht Querverweise, TOC-Eintraege und URLs
\usepackage[
  colorlinks = true,
  linkcolor  = tvpurple2,%linkblue,
  urlcolor   = tvpurple2,%linkblue,
  citecolor  = gray,
  pdftitle   = {TilVista Entwicklerdokumentation},
  pdfauthor  = {Ludwig Armin}
]{hyperref}

% HINWEIS: Diese Datei wird per \input{} eingebunden,
% nicht als eigenstaendiges Dokument kompiliert.

% ============================================================
%  macros.tex  --  Farben, Listings-Stile, tcolorbox-Umgebungen
%                  und Hilfsmakros fuer die TilVista-Dokumentation
%
%  Wird per % ============================================================
%  macros.tex  --  Farben, Listings-Stile, tcolorbox-Umgebungen
%                  und Hilfsmakros fuer die TilVista-Dokumentation
%
%  Wird per % ============================================================
%  macros.tex  --  Farben, Listings-Stile, tcolorbox-Umgebungen
%                  und Hilfsmakros fuer die TilVista-Dokumentation
%
%  Wird per \input{src/packages/macros} in tilvista_doku.tex
%  eingebunden, nach packages.tex und vor \begin{document}.
% ============================================================

% ============================================================
%  FARBEN
% ============================================================

% Icon Farbschema - TilVista
\definecolor{tvwhite}	{HTML}{d6d6d6}
\definecolor{tvpink1}	{HTML}{d8bfcd}
\definecolor{tvpink2}	{HTML}{daa9c8}
\definecolor{tvpink3}	{HTML}{dc8bc0}
\definecolor{tvpurple1}	{HTML}{a3add6}
\definecolor{tvpurple2}	{HTML}{646ec4}

% Hauptfarbe fuer C++-Schluesselwoerter und Links
\definecolor{cppblue}  {HTML}{003E7E}

% Qt-typisches Gruen (Signale, Qt-spezifische Klassen)
\definecolor{qtgreen}  {HTML}{2E7D32}

% Hintergrund fuer Code-Bloecke (sehr helles Grau-Blau)
\definecolor{codebg}   {HTML}{F4F6F9}

% Rahmenfarbe fuer Code-Bloecke
\definecolor{framerule}{HTML}{B0BEC5}

% Hinweis-Box (gelb)
\definecolor{calloutbg}{HTML}{FFF8E1}
\definecolor{calloutfr}{HTML}{F9A825}

% Warnungs-Box (rot)
\definecolor{warnbg}   {HTML}{FFEBEE}
\definecolor{warnfr}   {HTML}{C62828}

% Info-Box / Qt-Hintergrund (blau)
\definecolor{infobg}   {HTML}{E3F2FD}
\definecolor{infofr}   {HTML}{1565C0}

% Neue-Funktion-Box (gruen)
\definecolor{newbg}    {HTML}{E8F5E9}
\definecolor{newfr}    {HTML}{2E7D32}

% Link-Farbe (dunkles Marineblau)
\definecolor{linkblue} {HTML}{1A237E}

% CMake-Schluesselwoerter (Dunkelrot)
\definecolor{cmakered} {HTML}{8B1A1A}

% ============================================================
%  LISTINGS-STILE
% ============================================================

% -- Basis-Stil fuer C++ -------------------------------------
\lstdefinestyle{cpp}{
  language        = C++,
  basicstyle      = \ttfamily\small,
  % Schluesselwoerter fett + blau
  keywordstyle    = \bfseries\color{cppblue},
  % Kommentare kursiv + hellgrau
  commentstyle    = \itshape\color{gray!65},
  % Strings dunkelrot
  stringstyle     = \color{BrickRed},
  % Zeilennummern (klein, grau, links)
  numberstyle     = \tiny\color{gray!60},
  numbers         = left,
  numbersep       = 9pt,
  % Hintergrund und Rahmen
  backgroundcolor = \color{codebg},
  frame           = single,
  framerule       = 0.5pt,
  rulecolor       = \color{framerule},
  % Zeilenumbruch bei langen Zeilen
  breaklines      = true,
  breakatwhitespace = true,
  tabsize         = 4,
  showstringspaces = false,
  % WICHTIG: extendedchars=false verhindert, dass listings
  % versucht, UTF-8-Bytes selbst zu interpretieren.
  % Umlaute in lstlisting-Bloecken muessen als ASCII-Umschreibung
  % stehen (ae, oe, ue, ss). Im normalen Fliestext ausserhalb
  % von lstlisting funktionieren Umlaute dank inputenc normal.
  extendedchars   = false,
  % Zusaetzliche Qt- und C++-Schluesselwoerter
  morekeywords    = {signals,slots,emit,Q_OBJECT,override,
                     nullptr,auto,constexpr,explicit,
                     QString,QWidget,QTimer,QPixmap,QLabel,
                     QMainWindow,QObject,QThread,Qt,
                     bool,int,double,float,void},
  xleftmargin     = 2em,
  xrightmargin    = 0.5em,
  aboveskip       = 6pt,
  belowskip       = 6pt
}

% -- Stil fuer Header-Dateien (.h) ---------------------------
% Identisch mit cpp, aber mit leicht blaulichem Hintergrund,
% um Header optisch von .cpp-Dateien zu unterscheiden.
\lstdefinestyle{header}{
  style           = cpp,
  backgroundcolor = \color{blue!5}
}

% -- Stil fuer CMake-Dateien ---------------------------------
\lstdefinestyle{cmake}{
  language        = {},       % kein eingebautes CMake-Highlighting
  basicstyle      = \ttfamily\small,
  keywordstyle    = \bfseries\color{cmakered},
  commentstyle    = \itshape\color{gray!65},
  stringstyle     = \color{BrickRed},
  numberstyle     = \tiny\color{gray!60},
  numbers         = left,
  numbersep       = 9pt,
  backgroundcolor = \color{codebg},
  frame           = single,
  framerule       = 0.5pt,
  rulecolor       = \color{framerule},
  breaklines      = true,
  tabsize         = 2,
  showstringspaces = false,
  extendedchars   = false,
  morekeywords    = {cmake_minimum_required,project,set,
                     find_package,add_executable,
                     target_link_libraries,target_compile_definitions,
                     target_include_directories,
                     set_target_properties,
                     add_custom_command,add_compile_definitions,
                     function,endfunction,if,else,endif,
                     find_program,message,include,install},
  xleftmargin     = 2em,
  xrightmargin    = 0.5em,
  aboveskip       = 6pt,
  belowskip       = 6pt
}

% Standard-Stil fuer alle \begin{lstlisting} ohne style=...
\lstset{style=cpp}

% ============================================================
%  tcolorbox-UMGEBUNGEN (farbige Infokaesten)
% ============================================================

% -- Hinweis (gelb) ------------------------------------------
% Fuer allgemeine Tipps und Erklaerungen
\newtcolorbox{hinweis}[1][]{
  enhanced, breakable,
  colback    = calloutbg,
  colframe   = calloutfr,
  fonttitle  = \bfseries\small,
  title      = {Hinweis},
  left=6pt, right=6pt, top=4pt, bottom=4pt,
  #1
}

% -- Warnung (rot) -------------------------------------------
% Fuer Fallstricke, undefiniertes Verhalten, Deadlock-Gefahren
\newtcolorbox{warnung}[1][]{
  enhanced, breakable,
  colback    = warnbg,
  colframe   = warnfr,
  fonttitle  = \bfseries\small,
  title      = {Achtung},
  left=6pt, right=6pt, top=4pt, bottom=4pt,
  #1
}

% -- Qt-Hintergrund (blau) -----------------------------------
% Fuer Qt-spezifische Konzepterklarungen mit Doku-Links
\newtcolorbox{qtinfo}[1][]{
  enhanced, breakable,
  colback    = infobg,
  colframe   = infofr,
  fonttitle  = \bfseries\small,
  title      = {Qt-Hintergrund},
  left=6pt, right=6pt, top=4pt, bottom=4pt,
  #1
}

% -- Neu in v0.5.42 (gruen) ----------------------------------
% Hebt Aenderungen der aktuellen Version hervor
\newtcolorbox{neuin}[1][]{
  enhanced, breakable,
  colback    = newbg,
  colframe   = newfr,
  fonttitle  = \bfseries\small,
  title      = {Neu in v0.5.42},
  left=6pt, right=6pt, top=4pt, bottom=4pt,
  #1
}

% ============================================================
%  HILFSMAKROS (Semantische Auszeichnung)
% ============================================================

% Dateinamen: monospace, z.B. \datei{mainwindow.h}
\newcommand{\datei}[1]{\texttt{#1}}

% Klassennamen: monospace + fett, z.B. \klasse{MainWindow}
\newcommand{\klasse}[1]{\texttt{\textbf{#1}}}

% Methoden mit Klammern: \methode{buildUi} --> buildUi()
\newcommand{\methode}[1]{\texttt{#1()}}

% Member-Variablen: \member{m\_dirBar}
\newcommand{\member}[1]{\texttt{#1}}

% Signale (Qt-gruen): \signal{dirLoaded}
\newcommand{\signal}[1]{\texttt{\color{qtgreen}#1}}

% Slots (blau): \slot{onDirChanged}
\newcommand{\slot}[1]{\texttt{\color{cppblue}#1}}

% Praeprozessor-Makros (dunkelrot): \makro{TV\_NO\_MULTIMEDIA}
\newcommand{\makro}[1]{\texttt{\color{BrickRed}#1}}

% CMake-Befehle: \cmake{find\_package}
\newcommand{\cmake}[1]{\texttt{\color{cmakered}#1}}

% Namespace-Praefix: \ns{TV} --> TV::
\newcommand{\ns}[1]{\texttt{#1::}}

% Versionsnummer hervorgehoben: \version{0.5.42}
\newcommand{\version}[1]{\textbf{v#1}}

% ============================================================
%  SEITENSTIL (fancyhdr)
% ============================================================

\pagestyle{fancy}
\fancyhf{}  % alle Felder leeren

% Linke Seite (gerade Seitenzahl): Icon + Kapitelname
\fancyhead[LE]{%
  \raisebox{-2pt}{\includegraphics[height=11pt]{graphics/icon}}%
  \quad\small\leftmark%
}
% Rechte Seite (ungerade Seitenzahl): Abschnittsname
\fancyhead[RO]{\small\rightmark}
% Fusszeile zentriert: Seitenzahl
\fancyfoot[C]{\small\thepage}

\renewcommand{\headrulewidth}{0.4pt}
\renewcommand{\footrulewidth}{0pt}
 in tilvista_doku.tex
%  eingebunden, nach packages.tex und vor \begin{document}.
% ============================================================

% ============================================================
%  FARBEN
% ============================================================

% Icon Farbschema - TilVista
\definecolor{tvwhite}	{HTML}{d6d6d6}
\definecolor{tvpink1}	{HTML}{d8bfcd}
\definecolor{tvpink2}	{HTML}{daa9c8}
\definecolor{tvpink3}	{HTML}{dc8bc0}
\definecolor{tvpurple1}	{HTML}{a3add6}
\definecolor{tvpurple2}	{HTML}{646ec4}

% Hauptfarbe fuer C++-Schluesselwoerter und Links
\definecolor{cppblue}  {HTML}{003E7E}

% Qt-typisches Gruen (Signale, Qt-spezifische Klassen)
\definecolor{qtgreen}  {HTML}{2E7D32}

% Hintergrund fuer Code-Bloecke (sehr helles Grau-Blau)
\definecolor{codebg}   {HTML}{F4F6F9}

% Rahmenfarbe fuer Code-Bloecke
\definecolor{framerule}{HTML}{B0BEC5}

% Hinweis-Box (gelb)
\definecolor{calloutbg}{HTML}{FFF8E1}
\definecolor{calloutfr}{HTML}{F9A825}

% Warnungs-Box (rot)
\definecolor{warnbg}   {HTML}{FFEBEE}
\definecolor{warnfr}   {HTML}{C62828}

% Info-Box / Qt-Hintergrund (blau)
\definecolor{infobg}   {HTML}{E3F2FD}
\definecolor{infofr}   {HTML}{1565C0}

% Neue-Funktion-Box (gruen)
\definecolor{newbg}    {HTML}{E8F5E9}
\definecolor{newfr}    {HTML}{2E7D32}

% Link-Farbe (dunkles Marineblau)
\definecolor{linkblue} {HTML}{1A237E}

% CMake-Schluesselwoerter (Dunkelrot)
\definecolor{cmakered} {HTML}{8B1A1A}

% ============================================================
%  LISTINGS-STILE
% ============================================================

% -- Basis-Stil fuer C++ -------------------------------------
\lstdefinestyle{cpp}{
  language        = C++,
  basicstyle      = \ttfamily\small,
  % Schluesselwoerter fett + blau
  keywordstyle    = \bfseries\color{cppblue},
  % Kommentare kursiv + hellgrau
  commentstyle    = \itshape\color{gray!65},
  % Strings dunkelrot
  stringstyle     = \color{BrickRed},
  % Zeilennummern (klein, grau, links)
  numberstyle     = \tiny\color{gray!60},
  numbers         = left,
  numbersep       = 9pt,
  % Hintergrund und Rahmen
  backgroundcolor = \color{codebg},
  frame           = single,
  framerule       = 0.5pt,
  rulecolor       = \color{framerule},
  % Zeilenumbruch bei langen Zeilen
  breaklines      = true,
  breakatwhitespace = true,
  tabsize         = 4,
  showstringspaces = false,
  % WICHTIG: extendedchars=false verhindert, dass listings
  % versucht, UTF-8-Bytes selbst zu interpretieren.
  % Umlaute in lstlisting-Bloecken muessen als ASCII-Umschreibung
  % stehen (ae, oe, ue, ss). Im normalen Fliestext ausserhalb
  % von lstlisting funktionieren Umlaute dank inputenc normal.
  extendedchars   = false,
  % Zusaetzliche Qt- und C++-Schluesselwoerter
  morekeywords    = {signals,slots,emit,Q_OBJECT,override,
                     nullptr,auto,constexpr,explicit,
                     QString,QWidget,QTimer,QPixmap,QLabel,
                     QMainWindow,QObject,QThread,Qt,
                     bool,int,double,float,void},
  xleftmargin     = 2em,
  xrightmargin    = 0.5em,
  aboveskip       = 6pt,
  belowskip       = 6pt
}

% -- Stil fuer Header-Dateien (.h) ---------------------------
% Identisch mit cpp, aber mit leicht blaulichem Hintergrund,
% um Header optisch von .cpp-Dateien zu unterscheiden.
\lstdefinestyle{header}{
  style           = cpp,
  backgroundcolor = \color{blue!5}
}

% -- Stil fuer CMake-Dateien ---------------------------------
\lstdefinestyle{cmake}{
  language        = {},       % kein eingebautes CMake-Highlighting
  basicstyle      = \ttfamily\small,
  keywordstyle    = \bfseries\color{cmakered},
  commentstyle    = \itshape\color{gray!65},
  stringstyle     = \color{BrickRed},
  numberstyle     = \tiny\color{gray!60},
  numbers         = left,
  numbersep       = 9pt,
  backgroundcolor = \color{codebg},
  frame           = single,
  framerule       = 0.5pt,
  rulecolor       = \color{framerule},
  breaklines      = true,
  tabsize         = 2,
  showstringspaces = false,
  extendedchars   = false,
  morekeywords    = {cmake_minimum_required,project,set,
                     find_package,add_executable,
                     target_link_libraries,target_compile_definitions,
                     target_include_directories,
                     set_target_properties,
                     add_custom_command,add_compile_definitions,
                     function,endfunction,if,else,endif,
                     find_program,message,include,install},
  xleftmargin     = 2em,
  xrightmargin    = 0.5em,
  aboveskip       = 6pt,
  belowskip       = 6pt
}

% Standard-Stil fuer alle \begin{lstlisting} ohne style=...
\lstset{style=cpp}

% ============================================================
%  tcolorbox-UMGEBUNGEN (farbige Infokaesten)
% ============================================================

% -- Hinweis (gelb) ------------------------------------------
% Fuer allgemeine Tipps und Erklaerungen
\newtcolorbox{hinweis}[1][]{
  enhanced, breakable,
  colback    = calloutbg,
  colframe   = calloutfr,
  fonttitle  = \bfseries\small,
  title      = {Hinweis},
  left=6pt, right=6pt, top=4pt, bottom=4pt,
  #1
}

% -- Warnung (rot) -------------------------------------------
% Fuer Fallstricke, undefiniertes Verhalten, Deadlock-Gefahren
\newtcolorbox{warnung}[1][]{
  enhanced, breakable,
  colback    = warnbg,
  colframe   = warnfr,
  fonttitle  = \bfseries\small,
  title      = {Achtung},
  left=6pt, right=6pt, top=4pt, bottom=4pt,
  #1
}

% -- Qt-Hintergrund (blau) -----------------------------------
% Fuer Qt-spezifische Konzepterklarungen mit Doku-Links
\newtcolorbox{qtinfo}[1][]{
  enhanced, breakable,
  colback    = infobg,
  colframe   = infofr,
  fonttitle  = \bfseries\small,
  title      = {Qt-Hintergrund},
  left=6pt, right=6pt, top=4pt, bottom=4pt,
  #1
}

% -- Neu in v0.5.42 (gruen) ----------------------------------
% Hebt Aenderungen der aktuellen Version hervor
\newtcolorbox{neuin}[1][]{
  enhanced, breakable,
  colback    = newbg,
  colframe   = newfr,
  fonttitle  = \bfseries\small,
  title      = {Neu in v0.5.42},
  left=6pt, right=6pt, top=4pt, bottom=4pt,
  #1
}

% ============================================================
%  HILFSMAKROS (Semantische Auszeichnung)
% ============================================================

% Dateinamen: monospace, z.B. \datei{mainwindow.h}
\newcommand{\datei}[1]{\texttt{#1}}

% Klassennamen: monospace + fett, z.B. \klasse{MainWindow}
\newcommand{\klasse}[1]{\texttt{\textbf{#1}}}

% Methoden mit Klammern: \methode{buildUi} --> buildUi()
\newcommand{\methode}[1]{\texttt{#1()}}

% Member-Variablen: \member{m\_dirBar}
\newcommand{\member}[1]{\texttt{#1}}

% Signale (Qt-gruen): \signal{dirLoaded}
\newcommand{\signal}[1]{\texttt{\color{qtgreen}#1}}

% Slots (blau): \slot{onDirChanged}
\newcommand{\slot}[1]{\texttt{\color{cppblue}#1}}

% Praeprozessor-Makros (dunkelrot): \makro{TV\_NO\_MULTIMEDIA}
\newcommand{\makro}[1]{\texttt{\color{BrickRed}#1}}

% CMake-Befehle: \cmake{find\_package}
\newcommand{\cmake}[1]{\texttt{\color{cmakered}#1}}

% Namespace-Praefix: \ns{TV} --> TV::
\newcommand{\ns}[1]{\texttt{#1::}}

% Versionsnummer hervorgehoben: \version{0.5.42}
\newcommand{\version}[1]{\textbf{v#1}}

% ============================================================
%  SEITENSTIL (fancyhdr)
% ============================================================

\pagestyle{fancy}
\fancyhf{}  % alle Felder leeren

% Linke Seite (gerade Seitenzahl): Icon + Kapitelname
\fancyhead[LE]{%
  \raisebox{-2pt}{\includegraphics[height=11pt]{graphics/icon}}%
  \quad\small\leftmark%
}
% Rechte Seite (ungerade Seitenzahl): Abschnittsname
\fancyhead[RO]{\small\rightmark}
% Fusszeile zentriert: Seitenzahl
\fancyfoot[C]{\small\thepage}

\renewcommand{\headrulewidth}{0.4pt}
\renewcommand{\footrulewidth}{0pt}
 in tilvista_doku.tex
%  eingebunden, nach packages.tex und vor \begin{document}.
% ============================================================

% ============================================================
%  FARBEN
% ============================================================

% Icon Farbschema - TilVista
\definecolor{tvwhite}	{HTML}{d6d6d6}
\definecolor{tvpink1}	{HTML}{d8bfcd}
\definecolor{tvpink2}	{HTML}{daa9c8}
\definecolor{tvpink3}	{HTML}{dc8bc0}
\definecolor{tvpurple1}	{HTML}{a3add6}
\definecolor{tvpurple2}	{HTML}{646ec4}

% Hauptfarbe fuer C++-Schluesselwoerter und Links
\definecolor{cppblue}  {HTML}{003E7E}

% Qt-typisches Gruen (Signale, Qt-spezifische Klassen)
\definecolor{qtgreen}  {HTML}{2E7D32}

% Hintergrund fuer Code-Bloecke (sehr helles Grau-Blau)
\definecolor{codebg}   {HTML}{F4F6F9}

% Rahmenfarbe fuer Code-Bloecke
\definecolor{framerule}{HTML}{B0BEC5}

% Hinweis-Box (gelb)
\definecolor{calloutbg}{HTML}{FFF8E1}
\definecolor{calloutfr}{HTML}{F9A825}

% Warnungs-Box (rot)
\definecolor{warnbg}   {HTML}{FFEBEE}
\definecolor{warnfr}   {HTML}{C62828}

% Info-Box / Qt-Hintergrund (blau)
\definecolor{infobg}   {HTML}{E3F2FD}
\definecolor{infofr}   {HTML}{1565C0}

% Neue-Funktion-Box (gruen)
\definecolor{newbg}    {HTML}{E8F5E9}
\definecolor{newfr}    {HTML}{2E7D32}

% Link-Farbe (dunkles Marineblau)
\definecolor{linkblue} {HTML}{1A237E}

% CMake-Schluesselwoerter (Dunkelrot)
\definecolor{cmakered} {HTML}{8B1A1A}

% ============================================================
%  LISTINGS-STILE
% ============================================================

% -- Basis-Stil fuer C++ -------------------------------------
\lstdefinestyle{cpp}{
  language        = C++,
  basicstyle      = \ttfamily\small,
  % Schluesselwoerter fett + blau
  keywordstyle    = \bfseries\color{cppblue},
  % Kommentare kursiv + hellgrau
  commentstyle    = \itshape\color{gray!65},
  % Strings dunkelrot
  stringstyle     = \color{BrickRed},
  % Zeilennummern (klein, grau, links)
  numberstyle     = \tiny\color{gray!60},
  numbers         = left,
  numbersep       = 9pt,
  % Hintergrund und Rahmen
  backgroundcolor = \color{codebg},
  frame           = single,
  framerule       = 0.5pt,
  rulecolor       = \color{framerule},
  % Zeilenumbruch bei langen Zeilen
  breaklines      = true,
  breakatwhitespace = true,
  tabsize         = 4,
  showstringspaces = false,
  % WICHTIG: extendedchars=false verhindert, dass listings
  % versucht, UTF-8-Bytes selbst zu interpretieren.
  % Umlaute in lstlisting-Bloecken muessen als ASCII-Umschreibung
  % stehen (ae, oe, ue, ss). Im normalen Fliestext ausserhalb
  % von lstlisting funktionieren Umlaute dank inputenc normal.
  extendedchars   = false,
  % Zusaetzliche Qt- und C++-Schluesselwoerter
  morekeywords    = {signals,slots,emit,Q_OBJECT,override,
                     nullptr,auto,constexpr,explicit,
                     QString,QWidget,QTimer,QPixmap,QLabel,
                     QMainWindow,QObject,QThread,Qt,
                     bool,int,double,float,void},
  xleftmargin     = 2em,
  xrightmargin    = 0.5em,
  aboveskip       = 6pt,
  belowskip       = 6pt
}

% -- Stil fuer Header-Dateien (.h) ---------------------------
% Identisch mit cpp, aber mit leicht blaulichem Hintergrund,
% um Header optisch von .cpp-Dateien zu unterscheiden.
\lstdefinestyle{header}{
  style           = cpp,
  backgroundcolor = \color{blue!5}
}

% -- Stil fuer CMake-Dateien ---------------------------------
\lstdefinestyle{cmake}{
  language        = {},       % kein eingebautes CMake-Highlighting
  basicstyle      = \ttfamily\small,
  keywordstyle    = \bfseries\color{cmakered},
  commentstyle    = \itshape\color{gray!65},
  stringstyle     = \color{BrickRed},
  numberstyle     = \tiny\color{gray!60},
  numbers         = left,
  numbersep       = 9pt,
  backgroundcolor = \color{codebg},
  frame           = single,
  framerule       = 0.5pt,
  rulecolor       = \color{framerule},
  breaklines      = true,
  tabsize         = 2,
  showstringspaces = false,
  extendedchars   = false,
  morekeywords    = {cmake_minimum_required,project,set,
                     find_package,add_executable,
                     target_link_libraries,target_compile_definitions,
                     target_include_directories,
                     set_target_properties,
                     add_custom_command,add_compile_definitions,
                     function,endfunction,if,else,endif,
                     find_program,message,include,install},
  xleftmargin     = 2em,
  xrightmargin    = 0.5em,
  aboveskip       = 6pt,
  belowskip       = 6pt
}

% Standard-Stil fuer alle \begin{lstlisting} ohne style=...
\lstset{style=cpp}

% ============================================================
%  tcolorbox-UMGEBUNGEN (farbige Infokaesten)
% ============================================================

% -- Hinweis (gelb) ------------------------------------------
% Fuer allgemeine Tipps und Erklaerungen
\newtcolorbox{hinweis}[1][]{
  enhanced, breakable,
  colback    = calloutbg,
  colframe   = calloutfr,
  fonttitle  = \bfseries\small,
  title      = {Hinweis},
  left=6pt, right=6pt, top=4pt, bottom=4pt,
  #1
}

% -- Warnung (rot) -------------------------------------------
% Fuer Fallstricke, undefiniertes Verhalten, Deadlock-Gefahren
\newtcolorbox{warnung}[1][]{
  enhanced, breakable,
  colback    = warnbg,
  colframe   = warnfr,
  fonttitle  = \bfseries\small,
  title      = {Achtung},
  left=6pt, right=6pt, top=4pt, bottom=4pt,
  #1
}

% -- Qt-Hintergrund (blau) -----------------------------------
% Fuer Qt-spezifische Konzepterklarungen mit Doku-Links
\newtcolorbox{qtinfo}[1][]{
  enhanced, breakable,
  colback    = infobg,
  colframe   = infofr,
  fonttitle  = \bfseries\small,
  title      = {Qt-Hintergrund},
  left=6pt, right=6pt, top=4pt, bottom=4pt,
  #1
}

% -- Neu in v0.5.42 (gruen) ----------------------------------
% Hebt Aenderungen der aktuellen Version hervor
\newtcolorbox{neuin}[1][]{
  enhanced, breakable,
  colback    = newbg,
  colframe   = newfr,
  fonttitle  = \bfseries\small,
  title      = {Neu in v0.5.42},
  left=6pt, right=6pt, top=4pt, bottom=4pt,
  #1
}

% ============================================================
%  HILFSMAKROS (Semantische Auszeichnung)
% ============================================================

% Dateinamen: monospace, z.B. \datei{mainwindow.h}
\newcommand{\datei}[1]{\texttt{#1}}

% Klassennamen: monospace + fett, z.B. \klasse{MainWindow}
\newcommand{\klasse}[1]{\texttt{\textbf{#1}}}

% Methoden mit Klammern: \methode{buildUi} --> buildUi()
\newcommand{\methode}[1]{\texttt{#1()}}

% Member-Variablen: \member{m\_dirBar}
\newcommand{\member}[1]{\texttt{#1}}

% Signale (Qt-gruen): \signal{dirLoaded}
\newcommand{\signal}[1]{\texttt{\color{qtgreen}#1}}

% Slots (blau): \slot{onDirChanged}
\newcommand{\slot}[1]{\texttt{\color{cppblue}#1}}

% Praeprozessor-Makros (dunkelrot): \makro{TV\_NO\_MULTIMEDIA}
\newcommand{\makro}[1]{\texttt{\color{BrickRed}#1}}

% CMake-Befehle: \cmake{find\_package}
\newcommand{\cmake}[1]{\texttt{\color{cmakered}#1}}

% Namespace-Praefix: \ns{TV} --> TV::
\newcommand{\ns}[1]{\texttt{#1::}}

% Versionsnummer hervorgehoben: \version{0.5.42}
\newcommand{\version}[1]{\textbf{v#1}}

% ============================================================
%  SEITENSTIL (fancyhdr)
% ============================================================

\pagestyle{fancy}
\fancyhf{}  % alle Felder leeren

% Linke Seite (gerade Seitenzahl): Icon + Kapitelname
\fancyhead[LE]{%
  \raisebox{-2pt}{\includegraphics[height=11pt]{graphics/icon}}%
  \quad\small\leftmark%
}
% Rechte Seite (ungerade Seitenzahl): Abschnittsname
\fancyhead[RO]{\small\rightmark}
% Fusszeile zentriert: Seitenzahl
\fancyfoot[C]{\small\thepage}

\renewcommand{\headrulewidth}{0.4pt}
\renewcommand{\footrulewidth}{0pt}


% ============================================================
\begin{document}
% ============================================================

% -- Titelseite ----------------------------------------------
\begin{titlepage}
  \centering
  \vspace*{2.5cm}

  % Icon oben auf der Titelseite
  \includegraphics[width=6cm]{graphics/icon}\\[1.5em]

  {\Huge\bfseries\color{tvpurple2} TilVista}\\[0.5em]\color{tvpink3}
  {\large Entwicklerdokumentation}\\[0.3em]
  {\normalsize\color{tvpink1}%\color{tvwhite}%\color{gray}
    Version 0.5.42 \textbullet{}
    Qt 6.11 \textbullet{}
    C++17 \textbullet{}
    CMake 3.20}\\[3cm]

  \begin{tcolorbox}[
    width=0.72\textwidth,
    colback=codebg, colframe=cppblue,
    boxrule=0.8pt, arc=4pt,
    left=10pt, right=10pt, top=8pt, bottom=8pt
  ]
  \centering\ttfamily\small
  project(TilVista VERSION 0.5.42 LANGUAGES CXX)\\
  setup\_qt\_target(\$\{PROJECT\_NAME\})
  \end{tcolorbox}

  \vfill
  \color{tvwhite}
  {\small Ludwig, Armin}\\[0.3em]
  {\small Juli 2026}
\end{titlepage}

% -- Inhaltsverzeichnis --------------------------------------
\tableofcontents
\clearpage

% ============================================================
%  Kapitel einbinden
% ============================================================

% input vs. include:
%   \input{}  -- fuegt den Inhalt direkt ein, kein Seitenumbruch
%   \include{} -- immer auf neuer Seite, erlaubt \includeonly
% Wir nutzen \input, da wir keine partielle Kompilierung brauchen.

% ============================================================
%  Kapitel 1: Allgemein -- TilVista und MainWindow
% ============================================================
\chapter{Allgemein: TilVista und das MainWindow}

\section{Projektueberblick}

TilVista ist eine Qt6-Desktop-Anwendung zur Verwaltung und
Anzeige von Mediendateien -- Bildern und Videos.
Der Name ist ein Kunstwort aus \emph{tilvelde}
(norw. zufaellig) und \emph{vista} (span./ital. Aussicht).
Version \version{0.5.42} ist die aktuelle Entwicklungsversion.

Das Programm gliedert sich in drei Haupt-Werkzeuge,
die als Tabs im Hauptfenster angeordnet sind:

\begin{itemize}[noitemsep]
  \item \textbf{AleaVue} -- Zufaellige Vollbild-Diashow aus
    einem Verzeichnis (Bilder, S-Taste zum Bookmarken)
  \item \textbf{SattumaPic} -- Zufaelliger Dateiaufruf mit
    Vorschau (beliebige Dateitypen)
  \item \textbf{MadoLudus} -- Konfigurations-Tab, oeffnet ein
    eigenstaendiges \klasse{MadoludusWindow} fuer Bilder und
    Videos mit Flip-Book-Modus (\textbf{neu ab v0.5.42})
\end{itemize}

Zwei persistente Datenbanken erganzen die Werkzeuge:
\begin{itemize}[noitemsep]
  \item \textbf{kaivo (DB1)} -- Verzeichnis-Datenbank mit
    gecachten Dateilisten (\datei{.dshow} / \datei{.catalogue})
  \item \textbf{shujuko (DB2)} -- Lesezeichen fuer einzelne
    Dateien, je kaivo-Eintrag eine eigene JSON-Datei
\end{itemize}

\section{Architekturueberblick}

Das Programm folgt einer flachen, signalgesteuerten Hierarchie.
\klasse{MainWindow} besitzt alle Tabs und gemeinsam genutzten
Komponenten. Tabs kommunizieren \emph{ausschliesslich} ueber
Signale und Slots mit \klasse{MainWindow} -- niemals direkt
untereinander.

\begin{center}
\ttfamily\small
\begin{tabular}{|c|}
\hline
\textbf{MainWindow (QMainWindow)} \\
\hline
DirBar -- globale Verzeichnisauswahl (oben) \\
\hline
QTabWidget \\
\quad AleaVueTab \textbar{} SattumaPicTab \textbar{}
MadoludusTab \textbar{} ShortcutsTab \\
\hline
StatusBar + SecretIndicator (unten) \\
\hline
\end{tabular}
\end{center}

\begin{neuin}
\textbf{v0.5.42-Refactor Madoludus:}
Der fruehre \klasse{MadolodosTab} war ein
Monolith (Konfiguration + Wiedergabe in einer Klasse).
Jetzt trennt \klasse{MadoludusTab} die Konfiguration vom
Wiedergabe-Fenster (\klasse{MadoludusWindow}), das beim
Start-Klick als eigenstaendiges \klasse{QMainWindow}
geoeffnet wird. Das Fenster ist rahmenlos, immer im
Vordergrund und per Maus verschiebbar.

Neu ebenfalls: \verb|closeEvent| in \klasse{MainWindow}
ruft \methode{flushPendingWrites} auf beiden Datenbanken auf,
damit laufende Hintergrund-Saves beim Beenden nicht
abgebrochen werden.
\end{neuin}

\klasse{ShujukoPanel} ist eine Besonderheit:
Es wird in \klasse{MainWindow} erstellt, aber als geteilter
Rohzeiger (shared, not owned) sowohl an
\klasse{AleaVueTab} als auch an \klasse{SattumaPicTab}
weitergereicht. Der Eigentuemer ist \klasse{MainWindow}.

\section{Dateien: mainwindow.h}

\begin{lstlisting}[style=header, caption={mainwindow.h -- vollstaendige Deklaration}]
#pragma once
#include <QMainWindow>
#include <QStringList>

// Forward-Deklarationen: kein #include noetig,
// solange nur Zeiger deklariert werden.
class DirBar;
class AleaVueTab;
class DirDatabasePanel;
class MadoludusTab;     // v0.5.42: war MadolodosTab
class ShortcutsTab;
class ShujukoPanel;
class SattumaPicTab;
class QCloseEvent;      // v0.5.42: neu (fuer closeEvent)
class QLabel;
class QShortcut;
class QTabWidget;

class MainWindow : public QMainWindow
{
    Q_OBJECT
public:
    explicit MainWindow(QWidget* parent = nullptr);

protected:
    // v0.5.42: wartet auf laufende Hintergrund-Saves
    void closeEvent(QCloseEvent* event) override;

private slots:
    void onDirChanged(const QString& path);
    void onDirFromDb(const QString& path);
    void onDb1FilesLoaded(const QString& path,
                          const QStringList& imageFiles,
                          const QStringList& allFiles);
    void onActiveEntryChanged(const QString& entryName,
                               const QString& sourceDir);
    void toggleSecretMode();

private:
    void updateSecretIndicator();

    DirBar*        m_dirBar        = nullptr;
    QTabWidget*    m_tabs          = nullptr;
    ShujukoPanel*  m_shujuko       = nullptr;
    AleaVueTab*    m_aleaVueTab    = nullptr;
    SattumaPicTab* m_sattumaPicTab = nullptr;
    MadoludusTab*  m_madoludusTab  = nullptr;
    ShortcutsTab*  m_shortcutsTab  = nullptr;

    QLabel*    m_secretIndicator = nullptr;
    QShortcut* m_secretShortcut  = nullptr;
    bool       m_secretMode      = false;
};
\end{lstlisting}

\subsubsection{C++-Analyse}

\paragraph{\#pragma once}
Ein nicht-standardisierter, aber universell unterstuetzter
Include-Guard. Statt
\verb|#ifndef MAINWINDOW_H / #define / #endif|
sorgt \verb|#pragma once| dafuer, dass die Datei pro
Kompiliereinheit nur einmal eingelesen wird -- kuerzer und
fehlerunanfaelliger.

\paragraph{Forward-Deklarationen}
Statt \verb|#include "aleavuetab.h"| reicht im Header
\verb|class AleaVueTab;| -- solange nur Zeiger auf den Typ
deklariert werden, muss der Compiler die Klasse noch nicht
vollstaendig kennen. Weniger Abhaengigkeiten, schnellere
Kompilierung (Pimpl-Idiom light).

\paragraph{Mitglieder-Initialisierung mit = nullptr}
In C++11 und spaeter kann man Pointer-Member direkt im
Klassenkoerper initialisieren. Das verhindert
uninitialisierte Zeiger (undefined behavior) und erspart
explizite Konstruktorinitialisierer.

\subsubsection{Qt-Analyse}

\paragraph{QMainWindow}
Die Basisklasse liefert Menueleiste, Toolbar-Bereich,
CentralWidget und Statusleiste -- vorgefertigt, kein
manueller Aufbau noetig. \klasse{MainWindow} erbt alles
durch \verb|: public QMainWindow|.

\paragraph{Q\_OBJECT}
Pflicht-Makro fuer jede Klasse mit Signalen oder Slots.
Der Meta Object Compiler (\texttt{moc}) generiert daraus
den C++-Code, der das Signal-Slot-System implementiert.
\cmake{AUTOMOC ON} in \datei{SetupQt.cmake} erledigt den
\texttt{moc}-Aufruf automatisch beim Build.

\paragraph{private slots}
Slots koennen jede Sichtbarkeit haben. \verb|private slots|
bedeutet: Die Slots koennen nur intern oder ueber
\verb|connect()| genutzt werden, nicht von aussen als
normale Methode aufgerufen.

\section{Dateien: mainwindow.cpp}

\begin{lstlisting}[caption={Konstruktor: Fenster aufbauen}]
MainWindow::MainWindow(QWidget* parent)
    : QMainWindow(parent)
{
    setWindowTitle(
        QString("TilVista  \u2013  %1")
        .arg(TV_APPVERSION_DISPLAY));
    setMinimumSize(900, 640);

    // Icon aus Qt-Ressourcen (resources.qrc)
    const QIcon appIcon(":/icons/tilvista.ico");
    if (!appIcon.isNull()) setWindowIcon(appIcon);

    // Zentrales Widget + vertikales Layout
    auto* central = new QWidget;
    setCentralWidget(central);
    auto* mv = new QVBoxLayout(central);
    mv->setContentsMargins(8,8,8,8);
    mv->setSpacing(4);

    m_dirBar = new DirBar;
    connect(m_dirBar, &DirBar::directoryChanged,
            this, &MainWindow::onDirChanged);
    mv->addWidget(m_dirBar);

    // ShujukoPanel: Eigentuemer MainWindow, geteilt mit Tabs
    m_shujuko = new ShujukoPanel(this);

    m_tabs = new QTabWidget;

    // Lambda gibt den aktuellen Verzeichnis-Zustand weiter
    m_aleaVueTab = new AleaVueTab(
        [this]{ return m_dirBar->directory(); }, m_shujuko);
    m_sattumaPicTab = new SattumaPicTab(
        [this]{ return m_dirBar->directory(); }, m_shujuko);
    m_madoludusTab  = new MadoludusTab(
        [this]{ return m_dirBar->directory(); });
    m_shortcutsTab  = new ShortcutsTab;

    // Signalkette: DB1 -> MainWindow -> alle betroffenen Tabs
    DirDatabasePanel* db1 = m_aleaVueTab->dbPanel();
    connect(m_aleaVueTab, &AleaVueTab::requestDirChange,
            this, &MainWindow::onDirFromDb);
    connect(db1, &DirDatabasePanel::dirLoaded,
            this, &MainWindow::onDb1FilesLoaded);
    connect(db1, &DirDatabasePanel::activeEntryChanged,
            this, &MainWindow::onActiveEntryChanged);
    connect(db1, &DirDatabasePanel::requestShujukoValidation,
            m_shujuko, &ShujukoPanel::validateFiles);

    m_tabs->addTab(m_aleaVueTab,    "\U0001f5bc  AleaVue");
    m_tabs->addTab(m_sattumaPicTab, "\U0001f3b2  SattumaPic");
    m_tabs->addTab(m_madoludusTab,  "\U0001f39e  MadoLudus");
    m_tabs->addTab(m_shortcutsTab,  "About");
    mv->addWidget(m_tabs);

    // Schloss-Symbol in der Statusleiste
    m_secretIndicator = new QLabel("\U0001f512");
    statusBar()->addPermanentWidget(m_secretIndicator);

    m_secretShortcut = new QShortcut(
        QKeySequence("Ctrl+Alt+F8"), this);
    m_secretShortcut->setContext(Qt::ApplicationShortcut);
    connect(m_secretShortcut, &QShortcut::activated,
            this, &MainWindow::toggleSecretMode);
}
\end{lstlisting}

\subsubsection{C++-Analyse: Lambda als Callback}

\begin{lstlisting}
m_aleaVueTab = new AleaVueTab(
    [this]{ return m_dirBar->directory(); }, m_shujuko);
\end{lstlisting}

Ein \textbf{Lambda-Ausdruck} wird als
\verb|std::function<QString()>| uebergeben:
\begin{itemize}[noitemsep]
  \item \verb|[this]| -- der Lambda-Closure haelt eine
    Kopie des \texttt{this}-Zeigers von \klasse{MainWindow}
  \item \verb|{ return m_dirBar->directory(); }| -- liefert
    beim Aufruf das aktuelle Verzeichnis aus der DirBar
\end{itemize}
Damit muss \klasse{AleaVueTab} kein \klasse{MainWindow}
kennen -- \emph{Dependency Inversion Principle}.

\subsubsection{v0.5.42 neu: closeEvent}

\begin{lstlisting}[caption={closeEvent: warten auf laufende Saves}]
void MainWindow::closeEvent(QCloseEvent* event) {
    // Blockiert kurz (Event-Loop pumpen), bis in-flight
    // Saves auf Disk abgeschlossen sind.
    // Details: DirDatabasePanel::flushPendingWrites().
    m_aleaVueTab->dbPanel()->flushPendingWrites();
    m_shujuko->flushPendingWrites();
    QMainWindow::closeEvent(event);  // Basisklasse!
}
\end{lstlisting}

Ohne diese Methode koennte ein "Save to DB"-Thread,
der genau beim Beenden des Programms lief, mitten im
Schreiben abgebrochen werden -- das JSON waere dann
korrupt.

\subsubsection{Qt-Analyse: Signalkette}

Wenn der Nutzer einen DB-Eintrag doppelklickt:

\begin{enumerate}[noitemsep]
  \item \klasse{DirDatabasePanel} emittiert
    \signal{dirLoaded(path, imageFiles, allFiles)}
  \item \slot{MainWindow::onDb1FilesLoaded} empfaengt
  \item Leitet \verb|path + allFiles| an
    \klasse{SattumaPicTab} weiter
  \item Leitet \verb|path + allFiles| an
    \klasse{MadoludusTab} weiter
\end{enumerate}

Alle Tabs bekommen die Daten -- ohne sich gegenseitig
zu kennen.

\section{Gemeinsam genutzte Hilfsdateien}

\subsection{pathutils.h / pathutils.cpp}

Alle anwendungsweiten Pfade, Dateinamen-Hilfsfunktionen
und plattformspezifische Calls sind im Namespace \ns{TV}
gebuendelt.

\begin{lstlisting}[style=header, caption={pathutils.h}]
namespace TV {
    QString scriptDir();    // Verzeichnis der .exe
    QString kaivoDir();     // .../optegnelser/kaivo/
    QString kaivoDbPath();  // .../kaivo/tilvista_dirs.json
    QString shujukoPath(const QString& kaivoEntryName);
    QString logDir();       // .../optegnelser/
    QString safeFilename(const QString& name);
    QString autoEntryName(const QString& dirPath);

    const QStringList& imageSuffixes(); // .jpg, .png, ...
    const QStringList& videoSuffixes(); // .mp4, .mkv, ...
    QStringList allMediaSuffixes();

    void openPath(const QString& path);
    void selectInExplorer(const QString& path);
    void preventSleep();  // Win32: SetThreadExecutionState
    void restoreSleep();
}
\end{lstlisting}

\paragraph{Statische lokale Variable -- Singleton-Initialisierung}
\begin{lstlisting}
const QStringList& imageSuffixes() {
    static const QStringList s = {
        ".jpg", ".jpeg", ".png", ".webp",
        ".heic", ".heif", ".avif", ".tif", ...
    };
    return s;
}
\end{lstlisting}
\verb|static| bei einer lokalen Variable: Initialisierung
genau einmal beim ersten Aufruf, danach im Speicher bis
Programmende. Rueckgabe als \verb|const QStringList&|
(Referenz) -- keine Kopie bei jedem Aufruf.

\paragraph{Raw-String-Literal R"(...)" und safeFilename}
\begin{lstlisting}
QString safeFilename(const QString& name) {
    static const QString bad = R"(\/:*?"<>|)";
    QString out;
    for (QChar c : name)
        out += bad.contains(c) ? QChar('_') : c;
    while (out.startsWith('.')) out.remove(0,1);
    return out.isEmpty()
        ? QStringLiteral("entry") : out;
}
\end{lstlisting}
\verb|R"(...)"| ist ein \textbf{Raw-String-Literal} (C++11):
Backslashes und Anfuehrungszeichen brauchen kein Escape.
\verb|QStringLiteral| erzeugt den String zur Compile-Zeit.

\paragraph{Plattformspezifischer Code mit Q\_OS\_WIN}
\begin{lstlisting}
void preventSleep() {
#ifdef Q_OS_WIN
    SetThreadExecutionState(
        ES_CONTINUOUS | ES_SYSTEM_REQUIRED);
#endif
}
\end{lstlisting}
Qt-eigene Praeprozessormakros (\verb|Q_OS_WIN|,
\verb|Q_OS_MAC|, \verb|Q_OS_LINUX|) machen den Code
portable: Auf Nicht-Windows-Systemen kompiliert die
Funktion als leere Huelle.

\subsection{config.h / config.cpp}

\klasse{Config} liest beim Programmstart einmalig
\datei{tilvista\_config.json} und stellt Einstellungen
als statische Klassenmitglieder bereit.

\begin{lstlisting}[style=header, caption={config.h}]
class Config {
public:
    static void     load();
    static QVariant get(const QString& key);
    static bool     getBool(const QString& key);
private:
    static QJsonObject s_data;
    static QJsonObject s_defaults;
};
\end{lstlisting}

\paragraph{Statische Membervariablen -- Klassenweites Singleton}
\verb|static QJsonObject s_data| gehoert der Klasse,
nicht einer Instanz. Es gibt genau ein \verb|s_data|
im gesamten Programm -- klassisches Singleton-Muster
ohne Overhead.

\begin{qtinfo}
\textbf{QJsonObject als Konfigurationsspeicher:}
Qt stellt \klasse{QJsonObject}, \klasse{QJsonArray} und
\klasse{QJsonDocument} fuer vollstaendiges JSON-Handling
bereit. Parsing:
\begin{lstlisting}
QJsonParseError err;
auto doc = QJsonDocument::fromJson(f.readAll(), &err);
if (err.error == QJsonParseError::NoError
    && doc.isObject())
    s_data = doc.object();
\end{lstlisting}
Dokumentation: \url{https://doc.qt.io/qt-6/qjsondocument.html}
\end{qtinfo}

% ============================================================
%  Kapitel 2: AleaVue -- Zufaellige Bildschirmpraesentation
% ============================================================
\chapter{Tab 1 -- AleaVue: Zufaellige Diashow}

% ============================================================
\section{Funktionsweise}
% ============================================================

AleaVue ist ein Vollbild-Zufalls-Diashow-Werkzeug.
Der Nutzer waehlt ein Verzeichnis, startet die Praesentation,
und TilVista zeigt alle enthaltenen Bilder in zufaelliger
Reihenfolge -- alle 4,5 Sekunden automatisch wechselnd.

\textbf{Ablauf Schritt fuer Schritt:}
\begin{enumerate}
  \item Nutzer setzt Verzeichnis in \klasse{DirBar}
        (Browse-Dialog oder Enter in der Pfad-Zeile)
  \item Klick auf \emph{Load from Directory} startet
        \klasse{ScanWorker} im Hintergrund-Thread --
        der GUI-Thread bleibt reaktiv
  \item Nach dem Scan stehen Bildpfade in
        \member{m\_imagePaths}, alle Pfade in
        \member{m\_allFiles}
  \item Alternativ: Doppelklick auf einen kaivo-Eintrag
        laedt gecachte Listen aus \datei{.dshow}/\datei{.catalogue}
        -- kein erneuter Scan noetig
  \item Klick auf \emph{Start AleaVue} oeffnet
        \klasse{SlideshowWindow} (Vollbild oder Fenster)
\end{enumerate}

\textbf{Tastaturkuerzel im Fenster:}
\begin{center}
\small
\begin{tabular}{ll}
\toprule
\textbf{Taste} & \textbf{Aktion} \\
\midrule
Pfeil rechts & Naechstes Bild (zufaellig) \\
Pfeil links  & Zurueck in der Anzeigehistorie \\
Leertaste    & Timer anhalten / fortsetzen \\
\textbf{S}   & Aktuelles Bild in shujuko (DB2) speichern \\
Esc          & Diashow schliessen, Sleep-Prevention zuruecksetzen \\
\bottomrule
\end{tabular}
\end{center}

% ============================================================
\section{File by File}
% ============================================================

\subsection{aleavuetab.h}

\begin{lstlisting}[style=header, caption={aleavuetab.h}]
#pragma once
#include <QStringList>
#include <QWidget>
#include <functional>

class QCheckBox;
class QLabel;
class QProgressBar;
class QPushButton;
class QThread;
class DirDatabasePanel;
class ShujukoPanel;
class SlideshowWindow;

class AleaVueTab : public QWidget
{
    Q_OBJECT
public:
    // std::function statt direkter MainWindow-Abhaengigkeit
    explicit AleaVueTab(std::function<QString()> getGlobalDir,
                        ShujukoPanel*            shujuko,
                        QWidget* parent = nullptr);

    void onDirectoryChanged(const QString& path);
    void loadFromDirectory();

    // Zugriff auf kaivo-Panel (fuer MainWindow-Signalverdrahtung)
    DirDatabasePanel* dbPanel() const { return m_dbPanel; }

signals:
    // Anfrage an MainWindow, DirBar auf diesen Pfad zu setzen
    void requestDirChange(const QString& path);

private slots:
    void onDbDirLoaded(const QString& path,
                       const QStringList& imageFiles,
                       const QStringList& allFiles);
    void onScanDone(bool ok,
                    QStringList imageFiles,
                    QStringList allFiles);
    void onScanThenOpen(bool ok,
                        QStringList imageFiles,
                        QStringList allFiles);
    void onStartSlideshow();

private:
    void buildUi();
    void scanDir(const QString& path);
    void openWindow(const QString& dir,
                    const QStringList& imagePaths);

    std::function<QString()> m_getGlobalDir;
    ShujukoPanel*            m_shujuko;  // shared, not owned

    QStringList m_imagePaths;
    QStringList m_allFiles;
    QString     m_lastScannedPath;

    QLabel*           m_lblStatus       = nullptr;
    QProgressBar*     m_pb              = nullptr;
    QPushButton*      m_btnStart        = nullptr;
    QCheckBox*        m_chkFullscreen   = nullptr;
    DirDatabasePanel* m_dbPanel         = nullptr;
    SlideshowWindow*  m_slideshowWindow = nullptr;
    QThread*          m_scanThread      = nullptr;
};
\end{lstlisting}

\subsubsection{C++-Analyse}

\paragraph{std::function als Dependency Inversion}
\verb|std::function<QString()>| ist ein Funktionsobjekt,
das jeden aufrufbaren Typ aufnehmen kann -- Lambda, Funktor,
freie Funktion. \klasse{AleaVueTab} braucht
\klasse{MainWindow} nicht zu kennen.
Das entspricht dem \emph{Dependency Inversion Principle}:
Abhaengigkeit von einer Abstraktion (Funktion), nicht von
einer konkreten Klasse.

\paragraph{ShujukoPanel* -- shared, not owned}
Der Kommentar ist semantisch wichtig.
\klasse{AleaVueTab} darf den Pointer \emph{nicht} deleten.
In modernem C++ wuerde man \verb|std::shared_ptr|
oder \verb|std::observer_ptr| verwenden; hier genuegt der
Rohzeiger, weil \klasse{MainWindow} als einziger
Eigentuemer klar definiert ist und laenger lebt als
\klasse{AleaVueTab}.

\paragraph{Zwei Scan-Slots}
\klasse{AleaVueTab} hat zwei verschiedene Slots fuer
Scan-Ergebnisse:
\begin{itemize}[noitemsep]
  \item \slot{onScanDone} -- Scan abgeschlossen, Dateiliste
    speichern, kein Fenster oeffnen (Hintergrund-Scan
    nach Verzeichniswechsel)
  \item \slot{onScanThenOpen} -- Scan abgeschlossen
    \emph{und} sofort Slideshow oeffnen (Nutzer hatte auf
    Start geklickt, bevor der Scan fertig war)
\end{itemize}

\subsubsection{Qt-Analyse}

\paragraph{DirDatabasePanel* dbPanel() const}
Getter mit \verb|const|-Qualifikation: Die Methode
veraendert kein Mitglied. \klasse{MainWindow} ruft ihn
ab, um Signal-Verbindungen zum DB-Panel herzustellen,
ohne dass \klasse{AleaVueTab} diese Verdrahtung selbst
kennen muss.

\paragraph{QProgressBar als Scan-Feedback}
\member{m\_pb} wird bei Scan-Start mit
\verb|setRange(0,100)|, \verb|setValue(0)| und
\verb|setVisible(true)| aktiviert. Fortschritt
kommt vom Signal \signal{progressChanged(int)} von
\klasse{ScanWorker} direkt per \verb|connect| an
\verb|m_pb->setValue| -- kein Code in \klasse{AleaVueTab}
noetig.

% ----------------------------------------------------------
\subsection{aleavuetab.cpp}
% ----------------------------------------------------------

\subsubsection{scanDir -- Worker-Thread starten}

\begin{lstlisting}[caption={scanDir: das Worker-Thread-Grundmuster in TilVista}]
void AleaVueTab::scanDir(const QString& path) {
    pbStart(m_pb);                 // Progressbar anzeigen
    m_btnStart->setEnabled(false); // Button sperren
    m_lastScannedPath = path;

    // Alten Thread entkoppeln (falls noch laufend)
    safeStop(m_scanThread, this);

    auto* worker = new ScanWorker(path);
    auto* thread = new QThread;
    m_scanThread = thread;

    // Worker in den neuen Thread verschieben
    worker->moveToThread(thread);

    // Verbindungen aufbauen
    connect(thread, &QThread::started,
            worker, &ScanWorker::run);
    connect(worker, &ScanWorker::progressChanged,
            m_pb,   &QProgressBar::setValue);
    connect(worker, &ScanWorker::resultReady,
            this,   &AleaVueTab::onScanDone);

    // Self-cleanup: Worker und Thread loeschen sich selbst
    connect(worker, &ScanWorker::resultReady,
            thread, &QThread::quit);
    connect(worker, &ScanWorker::resultReady,
            worker, &QObject::deleteLater);
    connect(thread, &QThread::finished,
            thread, &QObject::deleteLater);

    thread->start();
}
\end{lstlisting}

\begin{qtinfo}[title={Qt-Hintergrund: Worker-Object-Muster}]
Das \emph{Worker-Object-Muster} ist das Qt-empfohlene
Muster fuer Hintergrundarbeit ohne GUI-Blockade:
\begin{enumerate}[noitemsep]
  \item \klasse{QObject}-Subklasse (\klasse{ScanWorker})
    kapselt die Arbeit in \methode{run}
  \item \verb|moveToThread()| verschiebt den Worker
    in den Thread -- alle Slots laufen dann im Thread
  \item \signal{QThread::started} loest
    \slot{ScanWorker::run} aus
  \item Ergebnisse kommen per Signal im GUI-Thread an
    (queued connection, automatisch bei Thread-Grenze)
  \item \verb|deleteLater()| stellt thread-sicheres
    Aufraumen sicher -- nie \verb|delete| direkt aufrufen
\end{enumerate}
Niemals Qt-GUI-Methoden direkt aus einem Worker-Thread
aufrufen -- nur Signale senden!\\
Referenz: \url{https://doc.qt.io/qt-6/thread-basics.html}
\end{qtinfo}

\subsubsection{safeStop -- Schutz vor Stale-Signalen}

\begin{lstlisting}[caption={safeStop: verhindert Signale in ungueltige Slots}]
static void safeStop(QThread*& t, QObject* receiver) {
    if (!t) return;
    // Alle Verbindungen vom Thread zu diesem Objekt trennen
    t->disconnect(receiver);
    t = nullptr;
    // Der Thread laeuft weiter und raumt sich selbst auf,
    // sendet aber keine Signale mehr an "receiver".
}
\end{lstlisting}

Wenn ein neuer Scan gestartet wird, bevor der alte fertig ist,
koennten spaete Signale des alten Workers den neuen Zustand
korrumpieren. \verb|disconnect(receiver)| trennt alle
Signal-Verbindungen vom alten Thread zu \verb|this| --
ohne den Thread zu stoppen oder darauf zu warten.

\begin{warnung}[title={Achtung: Stale-Signal-Bug}]
Wuerde man \verb|safeStop| weglassen und einfach einen
zweiten Thread starten, koennte das so enden:
\begin{enumerate}[noitemsep]
  \item Scan 1 laeuft, braucht 3 Sekunden
  \item Nutzer waehlt neues Verzeichnis, Scan 2 startet
  \item Scan 2 endet nach 1 Sekunde: \verb|m_imagePaths|
    enthaelt Verzeichnis 2
  \item Scan 1 endet spaeter: \verb|m_imagePaths| wird
    mit Verzeichnis 1 ueberschrieben -- falscher Zustand!
\end{enumerate}
\end{warnung}

\subsubsection{openWindow -- SlideshowWindow erzeugen}

\begin{lstlisting}[caption={openWindow: Vollbild- oder Fenstermodus}]
void AleaVueTab::openWindow(const QString& directory,
                             const QStringList& imagePaths)
{
    const bool fullscreen = m_chkFullscreen->isChecked();
    m_slideshowWindow = new SlideshowWindow(
        directory,
        imagePaths,
        TV::logDir(),
        false,          // logErrors
        fullscreen,     // v0.5.42: Flag an Window weitergeben
        m_shujuko);

    fullscreen
        ? m_slideshowWindow->showFullScreen()
        : m_slideshowWindow->show();
}
\end{lstlisting}

\klasse{SlideshowWindow} erhaelt \member{m\_shujuko} als
Pointer, um beim S-Tastendruck direkt
\verb|m_shujuko->addBookmark(path)| aufrufen zu koennen.

% ----------------------------------------------------------
\subsection{slideshowwindow.h / slideshowwindow.cpp}
% ----------------------------------------------------------

\subsubsection{Klassendeklaration}

\begin{lstlisting}[style=header, caption={slideshowwindow.h -- Kern}]
class SlideshowWindow : public QMainWindow {
    Q_OBJECT
public:
    SlideshowWindow(const QString&     directory,
                    const QStringList& imagePaths,
                    const QString&     logPath,
                    bool               logErrors,
                    bool               fullscreen, // v0.5.42
                    ShujukoPanel*      shujuko,
                    QWidget*           parent = nullptr);
signals:
    void bookmarkAdded(const QString& path);

protected:
    void keyPressEvent(QKeyEvent* event) override;
    void resizeEvent(QResizeEvent* event) override;

private slots:
    void changeImage();

private:
    void displayImage(const QString& path);
    void goBack();
    void logError(const QString& path,
                  const QString& error);

    QString     m_directory;
    QStringList m_imagePaths;
    QString     m_logPath;
    bool        m_logErrors;
    bool        m_fullscreen = false;  // v0.5.42
    ShujukoPanel* m_shujuko;

    QLabel* m_label = nullptr;  // Anzeige-Widget
    QTimer* m_timer = nullptr;  // 4500ms Auto-Advance

    QStringList m_imageOrder;   // Anzeigeverlauf
    int         m_backtrack = 0;// < 0: in Verlauf zurueck
    QString     m_current;      // aktuell angezeigter Pfad

    int    m_showW = 1920, m_showH = 1080;
    double m_ratio = 16.0 / 9.0;
    static constexpr int kIntervalMs = 4500;
};
\end{lstlisting}

\subsubsection{C++-Analyse: Vollbild vs. Fenster (v0.5.42)}

\begin{lstlisting}[caption={Konstruktor: Vollbild- vs. Fensterlogik}]
SlideshowWindow::SlideshowWindow(...)
{
    TV::preventSleep();

    const QRect scr =
        QGuiApplication::primaryScreen()
        ->availableGeometry();

    if (m_fullscreen) {
        // Vollbild: Cursor verstecken (kein OS-Chrome)
        setCursor(Qt::BlankCursor);
        m_showW = scr.width();
        m_showH = scr.height();
    } else {
        // Fenster: Cursor sichtbar (Titelleiste braucht ihn)
        // Startet links oben bei halber Bildschirmgroesse
        m_showW = scr.width()  / 2;
        m_showH = scr.height() / 2;
        setGeometry(0, 0, m_showW, m_showH);
    }
    m_ratio = double(m_showW) / double(m_showH);
    // ...
    connect(this, &SlideshowWindow::bookmarkAdded,
            m_shujuko, &ShujukoPanel::addBookmark);

    m_timer->start();
    changeImage();
}
\end{lstlisting}

\verb|Qt::BlankCursor| versteckt den Mauszeiger vollstaendig --
nur im Vollbild sinnvoll, wo kein Fensterrahmen vorhanden ist.

\subsubsection{C++-Analyse: Zufallsreihenfolge und Rueckwaertsnavigation}

\begin{lstlisting}[caption={changeImage: Verlaufslogik mit m\_backtrack}]
void SlideshowWindow::changeImage() {
    if (m_imagePaths.isEmpty()) return;

    if (m_backtrack >= -1) {
        // Vorwaerts: neues zufaelliges Bild
        m_backtrack = 0;
        std::mt19937 rng(std::random_device{}());
        std::uniform_int_distribution<int>
            d(0, m_imagePaths.size()-1);
        const QString img = m_imagePaths.at(d(rng));
        m_imageOrder.append(img);  // in Verlauf speichern
        displayImage(img);
    } else {
        // Rueckwaerts: aus Verlauf wiederholen
        ++m_backtrack;
        const int idx =
            m_imageOrder.size() + m_backtrack;
        if (idx >= 0 && idx < m_imageOrder.size())
            displayImage(m_imageOrder.at(idx));
        else { m_backtrack = 0; changeImage(); }
    }
}
\end{lstlisting}

\member{m\_backtrack} ist ein negativer Offset in
\member{m\_imageOrder}: 0 = vorderster Eintrag,
$-1$ = eins zurueck, $-2$ = zwei zurueck usw.
Die Liste waechst vorwaerts, der Nutzer kann jederzeit
durch den Verlauf navigieren.

\begin{hinweis}
\textbf{std::mt19937 + std::random\_device:}
\verb|std::mt19937| (Mersenne Twister) ist ein
qualitativ hochwertiger Pseudozufallsgenerator --
deutlich besser als \verb|rand()|. Er wird mit einem
nicht-deterministischen Seed aus der Betriebssystem-Entropie
(\verb|std::random_device{}()|) initialisiert.
\verb|std::uniform_int_distribution<int>(a,b)| liefert
gleichverteilte Ganzzahlen in $[a, b]$ inklusive.\\
Referenz: \url{https://en.cppreference.com/w/cpp/numeric/random}
\end{hinweis}

\subsubsection{Qt-Analyse: QImageReader mit adaptiver Skalierung}

\begin{lstlisting}[caption={displayImage: effizientes Laden grosser Bilder}]
void SlideshowWindow::displayImage(const QString& path) {
    m_current = path;
    QImageReader reader(path);
    const QSize orig = reader.size();

    // Skalierung VOR dem Dekodieren setzen
    if (orig.isValid() &&
        (orig.width()  > m_showW ||
         orig.height() > m_showH))
    {
        const double s = qMin(
            double(m_showW) / orig.width(),
            double(m_showH) / orig.height());
        reader.setScaledSize(
            QSize(int(orig.width()  * s),
                  int(orig.height() * s)));
    }
    const QImage img = reader.read();
    if (img.isNull()) {
        logError(path, reader.errorString());
        changeImage();  // fehlerhaftes Bild ueberspringen
        return;
    }
    // Feinabstimmung per Seitenverhaeltnis
    QPixmap pm = QPixmap::fromImage(img);
    if (double(pm.width()) / double(pm.height()) > m_ratio)
        pm = pm.scaledToWidth(m_showW,
                              Qt::SmoothTransformation);
    else
        pm = pm.scaledToHeight(m_showH - 18,
                               Qt::SmoothTransformation);

    m_label->setPixmap(pm);
    m_timer->stop();
    m_timer->start(kIntervalMs); // Timer neu starten
    TV::preventSleep();
}
\end{lstlisting}

\begin{qtinfo}[title={Qt-Hintergrund: QImageReader vs. QPixmap(path)}]
\verb|QPixmap(path)| laedt das Bild vollstaendig in
seiner Originalaufloesung in den Speicher, bevor
skaliert wird. Bei einem 20-Megapixel-JPEG bedeutet
das: 20 MP $\times$ 4 Byte/Pixel = ca. 80 MB RAM nur
fuer ein Bild.\\
\verb|QImageReader::setScaledSize()| hingegen teilt dem
Dekoder mit, dass er \emph{beim Dekodieren} skalieren
soll -- JPEG-Dekoder koennen z.B. direkt auf 1/2, 1/4
oder 1/8 der Originalgroesse dekodieren. Das spart
erheblich Speicher und Zeit.\\
Referenz: \url{https://doc.qt.io/qt-6/qimagereader.html}
\end{qtinfo}

\subsubsection{Qt-Analyse: S-Taste und Bookmark-Signal}

\begin{lstlisting}[caption={keyPressEvent: S-Taste sendet Signal}]
void SlideshowWindow::keyPressEvent(QKeyEvent* event) {
    switch (event->key()) {
    case Qt::Key_Escape:
        setCursor(Qt::ArrowCursor);
        TV::restoreSleep();
        close();
        break;
    case Qt::Key_S:
        // emit traegt den Pfad, ShujukoPanel empfaengt ihn
        if (!m_current.isEmpty())
            emit bookmarkAdded(m_current);
        break;
    case Qt::Key_Space:
        m_timer->isActive()
            ? m_timer->stop()
            : m_timer->start(kIntervalMs);
        break;
    // ...
    default:
        QMainWindow::keyPressEvent(event); // Basisklasse!
    }
}
\end{lstlisting}

Im Konstruktor ist verbunden:
\begin{lstlisting}
connect(this,     &SlideshowWindow::bookmarkAdded,
        m_shujuko, &ShujukoPanel::addBookmark);
\end{lstlisting}

\klasse{SlideshowWindow} weiss nur, dass jemand auf das
Signal hoert -- die Implementierung von
\klasse{ShujukoPanel} ist egal. Loses Koppeln durch
Signale.

% ----------------------------------------------------------
\subsection{dirdatabasepanel.h / dirdatabasepanel.cpp (kaivo DB1)}
% ----------------------------------------------------------

\subsubsection{Aufgabe und JSON-Schema}

\klasse{DirDatabasePanel} speichert Verzeichniseintraege
mit gecachten Dateilisten. Die Datenbankdatei:

\begin{lstlisting}[language={},numbers=none,
  caption={tilvista\_dirs.json -- Struktur}]
{
  "entries": [
    {
      "name":        "Urlaub2025",
      "path":        "../../Bilder/Urlaub2025",
      "image_files": "Urlaub2025_img.dshow",
      "all_files":   "Urlaub2025_all.catalogue",
      "scanned_at":  "2025-08-01T14:23:00",
      "hidden":      false
    }
  ]
}
\end{lstlisting}

\datei{\_img.dshow} und \datei{\_all.catalogue} sind
zeilenweise Dateilisten im kaivo-Verzeichnis.
Beim Laden liest \klasse{CatalogueLoadWorker} diese
ein und liefert sie als \klasse{QStringList}.

\subsubsection{v0.5.42 Threading-Architektur}

\begin{lstlisting}[style=header,
  caption={Drei Thread-Pointer in DirDatabasePanel}]
QThread* m_thread    = nullptr; // save / index load
QThread* m_catThread = nullptr; // catalogue lesen
QThread* m_updThread = nullptr; // Verzeichnis neu scannen
\end{lstlisting}

Jede Operation laeuft in einem eigenen Thread.
Vor jeder neuen Operation wird der zugehoerige Thread
per \methode{safeStop} entkoppelt.

\begin{neuin}[title={Neu in v0.5.42: setBusy und flushPendingWrites}]
\textbf{setBusy(bool):} Waehrend ein Katalog-Load-Thread
laeuft, werden Save/Load/Delete/Update-Buttons und die
Eintragsliste deaktiviert. Das verhindert konkurrierende
Zugriffe auf \member{m\_db} (das \klasse{QJsonObject}).

\textbf{flushPendingWrites():} Beim App-Beenden pumpt
diese Methode die Event-Loop, bis der laufende Save-Thread
fertig ist:
\begin{lstlisting}
void DirDatabasePanel::flushPendingWrites() {
    if (!m_thread) return;
    QElapsedTimer t; t.start();
    while (m_thread && t.elapsed() < 3000)
        QCoreApplication::processEvents(
            QEventLoop::AllEvents, 50);
}
\end{lstlisting}
\verb|QThread::wait()| wuerde hier deadlocken:
Der Worker emittiert \signal{resultReady} als Queued
Connection zurueck in den Haupt-Thread. Dieser muss seine
Event-Loop \emph{laufen} lassen, damit das Signal
zugestellt und \verb|quit()| aufgerufen wird.
\verb|processEvents()| haelt genau genug der Event-Loop
am Leben, ohne den Thread zu blockieren.
\end{neuin}

\subsubsection{Signale von DirDatabasePanel}

\begin{lstlisting}[style=header]
signals:
    // Dateilisten geladen (aus Cache oder nach Scan)
    void dirLoaded(const QString& absPath,
                   const QStringList& imageFiles,
                   const QStringList& allFiles);
    // Aktiver Eintrag geaendert (fuer shujuko-Panel)
    void activeEntryChanged(const QString& entryName,
                             const QString& sourceDir);
    // shujuko soll Datei-Existenz pruefen
    void requestShujukoValidation();
    // Secret Mode hat sich geaendert
    void secretModeChanged(bool active);
\end{lstlisting}

\subsubsection{C++-Analyse: launchWorker -- DRY-Template-Helfer}

\begin{lstlisting}[caption={launchWorker: generische Thread-Startfunktion}]
template<typename Worker, typename Slot>
static QThread* launchWorker(Worker* worker,
                              QObject* receiver,
                              Slot slot)
{
    auto* thread = new QThread;
    worker->moveToThread(thread);
    QObject::connect(thread, &QThread::started,
                     worker, &Worker::run);
    QObject::connect(worker, &Worker::resultReady,
                     receiver, slot);
    QObject::connect(worker, &Worker::resultReady,
                     thread, &QThread::quit);
    QObject::connect(worker, &Worker::resultReady,
                     worker, &QObject::deleteLater);
    QObject::connect(thread, &QThread::finished,
                     thread, &QObject::deleteLater);
    thread->start();
    return thread;
}
\end{lstlisting}

\textbf{Warum ein Template?} Alle Worker-Klassen haben
denselben Thread-Start-Boilerplate. Das Template
(\verb|typename Worker, typename Slot|) macht die Funktion
fuer jeden Worker-Typ nutzbar ohne Code-Wiederholung
(DRY: Don't Repeat Yourself). Der Compiler instanziiert
fuer jeden Aufruf eine konkrete Version.

\subsubsection{C++-Analyse: rawName und regex}

\begin{lstlisting}[caption={rawName: Dekorations-Praefix entfernen}]
static QString rawName(const QString& displayName) {
    // Entfernt das "CIRCLE ""-Praefix von hidden entries
    static const QRegularExpression prefix("^[^ ]+ ");
    return QString(displayName).remove(prefix);
}
\end{lstlisting}

In der Eintragsliste werden versteckte Eintraege mit
dem Praefix \verb|"(kreis) "| versehen.
\verb|QRegularExpression| ist Qt's moderne Regex-Klasse
(PCRE2-basiert, Unicode-faehig). Die \verb|static|-Variable
stellt sicher, dass die Regex nur einmal kompiliert wird.

\subsubsection{Qt-Analyse: PendingCache-Struct}

\begin{lstlisting}[style=header,
  caption={PendingCache: aggregierter Zustand}]
struct PendingCache {
    QString path;
    QStringList imageFiles, allFiles;
    bool valid = false;
} m_pending;
\end{lstlisting}

Ein anonymes Struct direkt als Member-Variable --
das fasst zusammengehoerige Daten zusammen, ohne eine
benannte Klasse anlegen zu muessen. Das Flag \verb|valid|
zeigt an, ob gecachte Daten vorhanden sind, die noch
nicht in die DB geschrieben wurden.

% ============================================================
%  Kapitel 3: SattumaPic -- Zufaelliger Dateiaufruf
% ============================================================
\chapter{Tab 2 -- SattumaPic: Zufaelliger Dateiaufruf}

% ============================================================
\section{Funktionsweise}
% ============================================================

SattumaPic oeffnet zufaellige Dateien aus einem Verzeichnis.
Anders als AleaVue arbeitet es mit \emph{beliebigen}
Dateitypen -- nicht nur Bildern.
Eine Vorschau (Bild, Video oder Dateisystem-Icon) erscheint
im rechten Bereich. Das eingebettete \klasse{ShujukoPanel}
gibt Zugriff auf gebookmarkte Dateien.

\textbf{Ablauf:}
\begin{enumerate}[noitemsep]
  \item Verzeichnis setzen oder kaivo-Eintrag laden
  \item Klick auf \emph{Random File} -- zufaellige Datei
    aus \member{m\_allFiles}
  \item Optional: \emph{Auto-pick on dir select} waehlt
    sofort beim Setzen des Verzeichnisses
  \item Optional: \emph{Auto-open on random} oeffnet
    die Datei sofort mit dem Standard-Programm
  \item \emph{Select in File's Directory} oeffnet den
    Datei-Explorer und markiert die Datei
  \item \klasse{ShujukoPanel} links erlaubt direkten
    Zugriff auf gebookmarkte Dateien des aktiven
    kaivo-Eintrags
\end{enumerate}

% ============================================================
\section{File by File}
% ============================================================

\subsection{sattumapictab.h}

\begin{lstlisting}[style=header, caption={sattumapictab.h}]
#pragma once
#include <QStringList>
#include <QWidget>
#include <functional>

class QCheckBox;
class QLabel;
class QProgressBar;
class QPushButton;
class QThread;
class ShujukoPanel;
class VideoPreviewWidget;

class SattumaPicTab : public QWidget
{
    Q_OBJECT
public:
    explicit SattumaPicTab(
        std::function<QString()> getGlobalDir,
        ShujukoPanel*            shujuko,
        QWidget* parent = nullptr);

    // allFiles = {} als Default-Argument:
    // Aufruf ohne Cache oder mit Cache moeglich
    void onDirectoryChanged(const QString& path,
                            const QStringList& allFiles = {});

private slots:
    void onPickRandom();
    void onScanDone(bool ok,
                    QStringList imageFiles,
                    QStringList allFiles);
    void onOpenFile();
    void onSelectInDir();
    void onPreviewFromShujuko(const QString& path);

private:
    void buildUi();
    void scanAndPick(const QString& path);
    void displayFile(const QString& path);
    void updatePreview(const QString& path);

    std::function<QString()> m_getGlobalDir;
    ShujukoPanel*            m_shujuko;  // shared, not owned

    QString     m_randomFile;
    QStringList m_allFiles;

    QLabel*             m_lblFile     = nullptr;
    QProgressBar*       m_pb          = nullptr;
    QPushButton*        m_btnRandom   = nullptr;
    QPushButton*        m_btnOpenFile = nullptr;
    QPushButton*        m_btnOpenDir  = nullptr;
    QCheckBox*          m_chkAutoPick = nullptr;
    QCheckBox*          m_chkAutoOpen = nullptr;
    VideoPreviewWidget* m_preview     = nullptr;
    QLabel*             m_lblType     = nullptr;
    QThread*            m_scanThread  = nullptr;
};
\end{lstlisting}

\subsubsection{C++-Analyse: Default-Argument}

\begin{lstlisting}
void onDirectoryChanged(const QString& path,
                        const QStringList& allFiles = {});
\end{lstlisting}

\verb|= {}| ist ein Default-Argument -- wenn
\verb|allFiles| nicht uebergeben wird, erhaelt der
Parameter eine leere \klasse{QStringList}.
Zwei Aufrufer nutzen dieselbe Methode unterschiedlich:
\begin{itemize}[noitemsep]
  \item \klasse{MainWindow::onDirChanged} -- ohne allFiles
    (Verzeichnis manuell gesetzt, kein Cache vorhanden)
  \item \klasse{MainWindow::onDb1FilesLoaded} -- mit allFiles
    (aus DB geladen, Cache direkt uebernehmen, kein Scan)
\end{itemize}

% ----------------------------------------------------------
\subsection{sattumapictab.cpp}
% ----------------------------------------------------------

\subsubsection{onPickRandom -- C++11-Zufallsgenerator}

\begin{lstlisting}[caption={onPickRandom: Zufallsauswahl und Auto-Open}]
void SattumaPicTab::onPickRandom() {
    const QString dir = m_getGlobalDir();
    if (dir.isEmpty()) {
        QMessageBox::warning(this, "No Directory",
            "Please select a directory first.");
        return;
    }
    // Noch keine Dateiliste? Erst scannen, dann waehleneinen
    if (m_allFiles.isEmpty()) {
        scanAndPick(dir);
        return;
    }
    // C++11 Mersenne-Twister mit OS-Entropie-Seed
    std::mt19937 rng(std::random_device{}());
    std::uniform_int_distribution<int>
        dist(0, m_allFiles.size()-1);

    m_randomFile = m_allFiles.at(dist(rng));
    displayFile(m_randomFile);

    // Auto-Open: Datei sofort mit System-Programm oeffnen
    if (m_chkAutoOpen->isChecked()) onOpenFile();
}
\end{lstlisting}

\subsubsection{updatePreview -- Dateityp-Erkennung}

\begin{lstlisting}[caption={updatePreview: drei Zweige fuer Bild/Video/sonstige}]
void SattumaPicTab::updatePreview(const QString& path) {
    const QString ext =
        '.' + QFileInfo(path).suffix().toLower();

    // -- Bild? ------------------------------------------------
    if (TV::imageSuffixes().contains(ext)) {
        QImageReader reader(path);
        const QSize orig = reader.size();
        // Vorab-Skalierung fuer grosse Bilder
        if (orig.isValid() &&
            (orig.width() > 1024 || orig.height() > 1024))
        {
            const double s = qMin(1024.0/orig.width(),
                                  1024.0/orig.height());
            reader.setScaledSize(
                QSize(int(orig.width()  * s),
                      int(orig.height() * s)));
        }
        const QImage img = reader.read();
        if (!img.isNull()) {
            m_preview->showPixmap(
                QPixmap::fromImage(img),
                QString("Image  %1").arg(ext));
            return;
        }
    }

    // -- Video? -----------------------------------------------
    if (TV::videoSuffixes().contains(ext)) {
        m_preview->play(path);
        return;
    }

    // -- Beliebige Datei: System-Icon verwenden ---------------
    QFileIconProvider fip;
    const QIcon icon = fip.icon(QFileInfo(path));
    if (!icon.isNull())
        m_preview->showPixmap(
            icon.pixmap(QSize(128,128)),
            QString("File  %1").arg(ext));
    else
        m_preview->clearPreview();
}
\end{lstlisting}

\begin{qtinfo}[title={Qt-Hintergrund: QFileIconProvider}]
\klasse{QFileIconProvider} liefert das native
Betriebssystem-Icon fuer einen Dateityp -- exakt das Icon,
das der Windows Explorer, macOS Finder oder Linux-Dateimanager
anzeigen wuerde. TilVista muss die assoziierten Programme
nicht kennen. Das Icon kommt direkt vom OS.\\
Referenz: \url{https://doc.qt.io/qt-6/qfileiconprovider.html}
\end{qtinfo}

\subsubsection{Config-Integration: Auto-Pick und Auto-Open}

\begin{lstlisting}[caption={buildUi: Einstellungen aus Config laden}]
m_chkAutoPick = new QCheckBox("Auto-pick on dir select");
m_chkAutoPick->setChecked(
    Config::getBool("auto_pick_on_dir_select"));

m_chkAutoOpen = new QCheckBox("Auto-open on random");
m_chkAutoOpen->setChecked(
    Config::getBool("auto_open_on_random"));
\end{lstlisting}

Die Standardwerte kommen aus \klasse{Config::s\_defaults}:
\begin{itemize}[noitemsep]
  \item \verb|auto_pick_on_dir_select| -- Standard: \verb|true|
  \item \verb|auto_open_on_random| -- Standard: \verb|false|
\end{itemize}

Persistenz ueber Programmstarts hinweg erfolgt
\emph{ausschliesslich} ueber \datei{tilvista\_config.json}.
Es gibt keinen "Einstellungen speichern"-Button --
die Datei muss manuell bearbeitet werden.

% ----------------------------------------------------------
\subsection{videopreviewwidget.h / videopreviewwidget.cpp}
% ----------------------------------------------------------

\subsubsection{Aufbau: QStackedWidget mit zwei Seiten}

\klasse{VideoPreviewWidget} erbt von \klasse{QStackedWidget}
und verwaltet zwei Anzeigeseiten:

\begin{center}
\small
\begin{tabular}{cll}
\toprule
\textbf{Index} & \textbf{Widget} & \textbf{Verwendung} \\
\midrule
0 & \klasse{QVideoWidget} & Native Videowiedergabe \\
  &                       & (nur ohne \makro{TV\_NO\_MULTIMEDIA}) \\
1 & \klasse{PreviewLabel} & Einzelbilder und ffmpeg-Frames \\
\bottomrule
\end{tabular}
\end{center}

\begin{lstlisting}[style=header, caption={videopreviewwidget.h}]
class VideoPreviewWidget : public QStackedWidget {
    Q_OBJECT
public:
    explicit VideoPreviewWidget(QWidget* parent=nullptr);
    ~VideoPreviewWidget() override;

    void play(const QString& path);
    void showPixmap(const QPixmap& pm,
                    const QString& label = {});
    void clearPreview();

signals:
    void statusText(const QString& text);

private slots:
    void onFramesReady(bool ok, const QStringList& paths);
    void nextFrame();

private:
    void stop_();
    void cleanupTmp();

    QMediaPlayer* m_player    = nullptr;
    QAudioOutput* m_audio     = nullptr;
    QVideoWidget* m_videoW    = nullptr;
    PreviewLabel* m_frameLabel;
    QTimer*       m_cycleTimer;

    QStringList m_framePaths;
    int         m_frameIdx = 0;
    QString     m_tmpDir;
    QThread*    m_vidThread = nullptr;

    // 10 Sekunden pro Frame im Vorschau-Modus
    static constexpr int kCycleMs = 10000;
};
\end{lstlisting}

\subsubsection{C++-Analyse: constexpr statt \#define}

\begin{lstlisting}
static constexpr int kCycleMs = 10000;
\end{lstlisting}

\verb|constexpr| ist die moderne C++-Alternative zu
\verb|#define CYCLE_MS 10000|:
\begin{itemize}[noitemsep]
  \item \textbf{Typsicher}: Der Compiler prueft den Typ
    (\verb|int|), \verb|#define| ist text-Substitution
  \item \textbf{Scoped}: \verb|kCycleMs| gehoert zur
    Klasse, kein globaler Namensraumverschmutz
  \item \textbf{Debuggable}: Debugger kennen den Namen,
    bei \verb|#define| sieht man nur die Zahl
\end{itemize}

\subsubsection{play() -- nativer oder ffmpeg-Fallback}

\begin{lstlisting}[caption={play: zwei Codepfade je nach Multimedia-Verfuegbarkeit}]
void VideoPreviewWidget::play(const QString& path) {
    stop_();
    cleanupTmp();

#ifndef TV_NO_MULTIMEDIA
    // Nativer Pfad: QMediaPlayer stumm, Endlos-Loop
    if (m_player) {
        setCurrentIndex(0);  // QVideoWidget-Seite
        m_player->setSource(QUrl::fromLocalFile(path));
        m_player->play();
        emit statusText("Video  (QMediaPlayer -- muted)");
        return;
    }
#endif

    // Fallback: ffmpeg-Frames extrahieren und zyklisch anzeigen
    setCurrentIndex(1);  // PreviewLabel-Seite
    emit statusText("Video -- extracting frames...");

    auto* w = new VideoFramesWorker(path);
    m_tmpDir = w->tmpDir();
    auto* t  = new QThread;
    m_vidThread = t;
    w->moveToThread(t);
    connect(t, &QThread::started,
            w, &VideoFramesWorker::run);
    connect(w, &VideoFramesWorker::resultReady,
            this, &VideoPreviewWidget::onFramesReady);
    connect(w, &VideoFramesWorker::resultReady,
            t, &QThread::quit);
    connect(w, &VideoFramesWorker::resultReady,
            w, &QObject::deleteLater);
    connect(t, &QThread::finished,
            t, &QObject::deleteLater);
    t->start();
}
\end{lstlisting}

\subsubsection{Qt-Analyse: stop\_() und cleanupTmp()}

\begin{lstlisting}[caption={stop\_ und cleanupTmp: aufraeuemen vor neuem Inhalt}]
void VideoPreviewWidget::stop_() {
    m_cycleTimer->stop();
#ifndef TV_NO_MULTIMEDIA
    if (m_player) m_player->stop();
#endif
    if (m_vidThread) {
        m_vidThread->quit();
        m_vidThread = nullptr;
        m_vidWorker = nullptr;
    }
}

void VideoPreviewWidget::cleanupTmp() {
    if (!m_tmpDir.isEmpty()) {
        // Gesamtes temp-Verzeichnis rekursiv loeschen
        QDir(m_tmpDir).removeRecursively();
        m_tmpDir.clear();
    }
    m_framePaths.clear();
    m_frameIdx = 0;
}
\end{lstlisting}

\verb|QDir::removeRecursively()| loescht ein Verzeichnis
mit allem Inhalt -- das Aequivalent zu \texttt{rm -rf}.
Wichtig: Immer \methode{stop\_} und dann \methode{cleanupTmp}
aufrufen, bevor neue Inhalte geladen werden, sonst wuerde
ein laufender Frame-Timer auf geloeschte Dateien zeigen.

% ----------------------------------------------------------
\subsection{previewlabel.h / previewlabel.cpp}
% ----------------------------------------------------------

\klasse{PreviewLabel} erbt von \klasse{QLabel} und
ueberschreibt \methode{resizeEvent}, um das Bild bei
Groessenaenderung automatisch neu zu skalieren.

\begin{lstlisting}[style=header, caption={previewlabel.h}]
class PreviewLabel : public QLabel {
    Q_OBJECT
public:
    explicit PreviewLabel(QWidget* parent = nullptr);
    void setSourcePixmap(const QPixmap& pm);
    void clearPreview();
protected:
    void resizeEvent(QResizeEvent* e) override;
private:
    void repaintScaled();
    QPixmap m_source;  // Originalbild (ungekuerzt)
};
\end{lstlisting}

\begin{lstlisting}[caption={previewlabel.cpp -- Kern}]
void PreviewLabel::setSourcePixmap(const QPixmap& pm) {
    m_source = pm;
    repaintScaled();
}
void PreviewLabel::resizeEvent(QResizeEvent* e) {
    QLabel::resizeEvent(e);  // Basisklasse aufrufen!
    repaintScaled();
}
void PreviewLabel::repaintScaled() {
    if (m_source.isNull()) return;
    setPixmap(m_source.scaled(
        size() - QSize(8,8),           // 4px Rand ringsum
        Qt::KeepAspectRatio,
        Qt::SmoothTransformation));
}
\end{lstlisting}

\textbf{Das Muster}: \member{m\_source} haelt das
Originalbild. Bei jedem \verb|resizeEvent| wird eine
neu skalierte Kopie erzeugt und als Pixmap gesetzt.
So sieht die Vorschau immer scharf aus, egal in welcher
Fenstergroesse.

\begin{hinweis}
Beim Ueberschreiben von Event-Handlern muss die Methode
der Basisklasse aufgerufen werden
(\verb|QLabel::resizeEvent(e)|), damit die Basisklasse
ihr eigenes Layout-Update durchfuehren kann. Wird dies
vergessen, koennen Layoutfehler oder Rendering-Artefakte
entstehen.
\end{hinweis}

% ----------------------------------------------------------
\subsection{shujukopanel.h / shujukopanel.cpp (DB2)}
% ----------------------------------------------------------

\klasse{ShujukoPanel} (DB2) ist das Lesezeichen-System.
Es erscheint eingebettet im linken Bereich von
\klasse{SattumaPicTab}.

\subsubsection{JSON-Schema}

\begin{lstlisting}[language={},numbers=none,
  caption={<entryname>\_shujuko.json}]
{
  "kaivo_entry": "Urlaub2025",
  "source_dir":  "E:/Bilder/Urlaub2025",
  "files": [
    {
      "path":     "E:/Bilder/.../IMG_001.jpg",
      "added_at": "2025-08-01T15:30:00",
      "missing":  false
    }
  ]
}
\end{lstlisting}

Jeder kaivo-Eintrag hat eine eigene shujuko-JSON-Datei.
\member{missing} wird von \methode{validateFiles}
aktualisiert.

\subsubsection{Klassendeklaration}

\begin{lstlisting}[style=header, caption={shujukopanel.h -- Kern}]
class ShujukoPanel : public QWidget {
    Q_OBJECT
public:
    explicit ShujukoPanel(QWidget* parent = nullptr);

    void setActiveKaivoEntry(const QString& entryName,
                              const QString& sourceDir);
    void addBookmark(const QString& filePath);
    void validateFiles();
    void setSecretMode(bool on);
    QString currentJsonPath() const;

    // v0.5.42: blockiert kurz bis laufender Save fertig
    void flushPendingWrites();

signals:
    void fileSelected(const QString& path);

private:
    // ...Slots und Private-Methoden...
    QString     m_entryName;
    QString     m_sourceDir;
    QString     m_jsonPath;
    QJsonObject m_data;

    QThread* m_thread = nullptr;
    bool     m_saving = false;  // Guard: kein Doppel-Save
};
\end{lstlisting}

\subsubsection{C++-Analyse: Bool-Guard gegen Doppel-Save}

\begin{lstlisting}[caption={saveToDisk: m\_saving als einfacher Guard}]
void ShujukoPanel::saveToDisk() {
    if (m_jsonPath.isEmpty() || m_saving) return;
    m_saving = true;  // Guard setzen

    safeStopThread();
    auto* worker = new DBSaveWorker(m_jsonPath, m_data);
    auto* thread = new QThread;
    m_thread = thread;
    // ...Verbindungen...
    thread->start();
}

void ShujukoPanel::onSaveDone(bool ok) {
    m_saving = false;  // Guard zuruecksetzen
    m_thread = nullptr;
}
\end{lstlisting}

Wenn der Nutzer mehrere S-Tasten-Druecke schnell hintereinander
ausfuehrt, wuerde ohne Guard jedes Mal ein neuer Save-Thread
gestartet. Mit \member{m\_saving} wird der zweite Aufruf
sofort abgewiesen -- der erste muss erst abgeschlossen sein.

\begin{warnung}[title={Achtung: Kein Mutex noetig?}]
Das funktioniert hier ohne \verb|std::mutex|, weil
\methode{saveToDisk} und \methode{onSaveDone} beide
im GUI-Thread laufen (verbunden als Auto-Connection
innerhalb desselben Threads). Es gibt keine echte
Race Condition -- der GUI-Thread ist single-threaded.
Wuerde \methode{saveToDisk} aus einem Worker-Thread
aufgerufen, waere ein Mutex zwingend noetig.
\end{warnung}

\subsubsection{Qt-Analyse: QJsonObject als Wert-Kopie}

\begin{lstlisting}
auto* worker = new DBSaveWorker(m_jsonPath, m_data);
\end{lstlisting}

\klasse{DBSaveWorker} haelt \verb|QJsonObject m_data|
als \emph{Kopie}. Das ist thread-sicher, weil Qt-Wertobjekte
(\klasse{QJsonObject}, \klasse{QString}, \klasse{QStringList})
\emph{Copy-on-Write} (implizites Sharing) verwenden:
Die Kopie ist guenstig (nur ein Referenzzaehler), und erst
wenn jemand schreibt, wird wirklich kopiert. Beide
Threads haben danach unabhaengige Kopien.

\subsubsection{validateFiles -- Datei-Existenz pruefen}

\begin{lstlisting}[caption={validateFiles: missing-Flag aktualisieren}]
void ShujukoPanel::validateFiles() {
    if (m_entryName.isEmpty()) return;
    QJsonArray files = m_data.value("files").toArray();
    bool changed = false;

    for (int i = 0; i < files.size(); ++i) {
        QJsonObject f = files[i].toObject();
        const bool exists =
            QFileInfo::exists(f.value("path").toString());
        // Zustand hat sich geaendert?
        if (f.value("missing").toBool() != !exists) {
            f["missing"] = !exists;
            files[i] = f;
            changed = true;
        }
    }
    if (changed) {
        m_data["files"] = files;
        refreshFileList();
        saveToDisk();
    }
}
\end{lstlisting}

\verb|QFileInfo::exists(path)| ist eine statische Methode --
kein Objekt noetig, ein einzelner Aufruf prueft die
Datei-Existenz. Fehlende Eintraege erscheinen in der Liste
mit Warn-Praefix und grauer Farbe.

% ============================================================
%  Kapitel 4: MadoLudus -- In-Window-Slideshow (v0.5.42)
% ============================================================
\chapter{Tab 3 -- MadoLudus: Konfiguration und Wiedergabefenster}

% ============================================================
\section{Funktionsweise}
% ============================================================

MadoLudus ist eine Diashow, die in einem \emph{eigenen,
rahmenlosen Fenster} (\klasse{MadoludusWindow}) laeuft.
Der Tab selbst ist nur die Konfigurationsoberflaeche -- das
Fenster oeffnet sich auf Knopfdruck.

\begin{neuin}[title={v0.5.42: Kompletter Refactor}]
Der fruehre \klasse{MadolodosTab} war ein Monolith:
Konfiguration und Wiedergabe in einer Klasse.
Ab v0.5.42 sind zwei Klassen getrennt:
\begin{itemize}[noitemsep]
  \item \klasse{MadoludusTab} -- Konfigurations-UI
    im Hauptfenster (bleibt immer sichtbar)
  \item \klasse{MadoludusWindow} -- Eigenstaendiges
    \klasse{QMainWindow}, rahmenlos, immer im Vordergrund,
    per Maus verschiebbar
\end{itemize}
Neu ebenfalls: Shuffle-Button (\textbf{R}-Taste) und
FF-Rate-Auswahl (1/sec, 1/2 sec, 1/N Frames) als Combobox.
\end{neuin}

\textbf{Betriebsmodi fuer Videos:}
\begin{center}
\small
\begin{tabular}{lll}
\toprule
\textbf{Modus} & \textbf{Bedingung} & \textbf{Mechanismus} \\
\midrule
Native Wiedergabe & FF-Mode aus, Multimedia vorhanden &
  \klasse{QMediaPlayer} + \klasse{QVideoWidget} \\
FF/Flip-Book & FF-Mode an (Standard) &
  ffmpeg-Frames via \klasse{VideoFramesWorker} \\
Fallback & kein Multimedia (\makro{TV\_NO\_MULTIMEDIA}) &
  immer ffmpeg-Frames \\
\bottomrule
\end{tabular}
\end{center}

\textbf{Tastaturkuerzel in MadoludusWindow:}
\begin{center}
\small
\begin{tabular}{ll}
\toprule
\textbf{Taste} & \textbf{Aktion} \\
\midrule
Space         & Play / Pause \\
Pfeil rechts  & Naechste Datei \\
Pfeil links   & Vorherige Datei (Verlauf-bewusst) \\
Shift+Rechts  & Video: +5 Sekunden (nur native) \\
Shift+Links   & Video: --5 Sekunden (nur native) \\
M             & Stummschaltung umschalten \\
F             & FF-Modus (Secret Mode zum Deaktivieren) \\
R             & Zufaellige / sequenzielle Reihenfolge \\
Esc           & Fenster schliessen \\
\bottomrule
\end{tabular}
\end{center}

% ============================================================
\section{File by File}
% ============================================================

\subsection{madoluduswindow.h -- MadoludusConfig-Struct}

\begin{lstlisting}[style=header,
  caption={madoluduswindow.h -- MadoludusConfig und Klassendeklaration}]
// Konfiguration, die MadoludusTab an MadoludusWindow uebergibt
struct MadoludusConfig {
    bool   showImages   = true;
    bool   showVideos   = true;
    bool   ffMode       = true;
    // Wert fuer Frame-Extraktion: Sekunden oder Anzahl
    double ffValue      = 10.0;
    // false = IntervalSeconds, true = FixedCount
    bool   ffFixedCount = false;
    bool   muted        = false;
    bool   randomMode   = false;
    bool   secretMode   = false;
};

class MadoludusWindow : public QMainWindow
{
    Q_OBJECT
public:
    explicit MadoludusWindow(
        const QStringList&    mediaFiles,
        const MadoludusConfig& cfg,
        QWidget*              parent = nullptr);
    ~MadoludusWindow() override;

signals:
    void windowClosed();
    void currentFileChanged(const QString& path);

protected:
    void keyPressEvent(QKeyEvent*  event) override;
    void resizeEvent(QResizeEvent* event) override;
    void mousePressEvent(QMouseEvent* event) override;
    void mouseMoveEvent(QMouseEvent*  event) override;
    void closeEvent(QCloseEvent*   event) override;

private slots:
    void onNextFile();
    void onPrevFile();
    void onPlayPause();
    void onMuteToggle();
    void onFFToggle();
    void onShuffleToggle();    // v0.5.42
    void onAutoAdvance();
    void onVideoEnded();
    void onFFFramesReady(bool ok, QStringList framePaths);
    void cycleFFFrame();

private:
    void buildUi();
    void loadFile(int index);
    void loadImage(const QString& path);
    void loadVideo(const QString& path);
    int  nextIndex() const;
    void startAutoAdvance();
    void stopAutoAdvance();
    void startFFExtraction(const QString& path);
    void stopFFExtraction();
    void updateButtonStates();
    void scaleCurrentImage();

    QStringList     m_mediaFiles;
    MadoludusConfig m_cfg;
    int             m_index        = -1;
    bool            m_paused       = false;
    bool            m_isVideo      = false;
    bool            m_usingFlipbook = false;

    // Navigations-Verlauf (Previous funktioniert auch im Zufallsmodus)
    QList<int> m_order;
    int        m_orderPos = -1;

    QWidget*     m_central   = nullptr;
    QLabel*      m_display   = nullptr;  // Bild / Flip-Book
    QLabel*      m_lblPath   = nullptr;  // Pfad unten
    MadoOverlay* m_overlay   = nullptr;  // linke Steuerleiste

    QTimer* m_advTimer     = nullptr;   // Auto-Advance
    QTimer* m_ffCycleTimer = nullptr;   // Frame-Wechsel

    QThread*    m_ffThread   = nullptr;
    QStringList m_ffFrames;
    int         m_ffFrameIdx = 0;
    QString     m_ffTmpDir;

#ifndef TV_NO_MULTIMEDIA
    QWidget*      m_videoContainer = nullptr;
    QMediaPlayer* m_player         = nullptr;
    QAudioOutput* m_audio          = nullptr;
    QVideoWidget* m_videoW         = nullptr;
#endif

    // Drag-to-move (rahmenlos)
    QPoint m_dragStart;
    bool   m_dragging = false;

    static constexpr int kImageDelay   = 5000; // ms
    static constexpr int kImageDelayFF = 2000;
    static constexpr int kFFCycleMs    = 600;
    static constexpr int kSeekStep     = 5000;
};
\end{lstlisting}

\subsubsection{C++-Analyse: Struct mit Default-Werten}

\begin{lstlisting}
struct MadoludusConfig {
    bool   ffMode  = true;
    double ffValue = 10.0;
    // ...
};
\end{lstlisting}

Ein \verb|struct| in C++ ist wie eine \verb|class|, aber
mit Standard-Sichtbarkeit \verb|public|. Die Member
werden direkt initialisiert (C++11). Das bedeutet:
\begin{lstlisting}
MadoludusConfig cfg;   // alle Member auf Default-Werte
cfg.muted = true;      // einzelne anpassen
\end{lstlisting}
Dieses Muster heisst \emph{Aggregate Initialization} --
kein Konstruktor noetig, solange kein
Verhalten gekapselt werden muss.

\subsubsection{C++-Analyse: Navigations-Verlauf}

\begin{lstlisting}[caption={nextIndex und Verlaufslogik}]
int MadoludusWindow::nextIndex() const {
    if (m_mediaFiles.size() <= 1)
        return qMax(0, m_index);
    if (!m_cfg.randomMode)
        return (m_index + 1) % m_mediaFiles.size();

    // Zufaelligen Index waehlen (nicht den aktuellen)
    std::mt19937 rng(std::random_device{}());
    std::uniform_int_distribution<int>
        dist(0, m_mediaFiles.size() - 1);
    int n;
    do { n = dist(rng); } while (n == m_index);
    return n;
}

void MadoludusWindow::onNextFile() {
    if (m_mediaFiles.isEmpty()) return;
    int next;
    if (m_orderPos < m_order.size() - 1) {
        // Verlauf vorwaerts (Nutzer hatte zurueck
        // navigiert und geht nun wieder vor)
        ++m_orderPos;
        next = m_order.at(m_orderPos);
    } else {
        // Neuen Index berechnen und in Verlauf aufnehmen
        next = nextIndex();
        m_order.append(next);
        m_orderPos = m_order.size() - 1;
    }
    loadFile(next);
}
\end{lstlisting}

Das Verlaufsystem mit \member{m\_order} (Liste der
angezeigten Indizes) und \member{m\_orderPos} (Position
darin) erlaubt beliebiges Vor- und Rueckwaertsnavigieren
-- auch im Zufallsmodus wiederholt man exakt die vorher
gezeigten Dateien.

\subsubsection{Qt-Analyse: Rahmenloses Fenster und Drag-to-Move}

\begin{lstlisting}[caption={Konstruktor: Fenster-Flags}]
setWindowFlags(Qt::Window
             | Qt::FramelessWindowHint
             | Qt::WindowStaysOnTopHint);
setAttribute(Qt::WA_DeleteOnClose);
setStyleSheet("background: black;");
\end{lstlisting}

\begin{itemize}[noitemsep]
  \item \verb|FramelessWindowHint| -- kein Titelbalken,
    kein OS-Rahmen; Fenster muss manuell bewegbar gemacht
    werden
  \item \verb|WindowStaysOnTopHint| -- immer im Vordergrund
    (auch vor anderen Anwendungen)
  \item \verb|WA_DeleteOnClose| -- das Objekt wird beim
    Schliessen automatisch geloescht (kein Speicherleck)
\end{itemize}

Drag-to-Move-Implementierung:
\begin{lstlisting}[caption={Mouse-Events fuer Drag-to-Move}]
void MadoludusWindow::mousePressEvent(QMouseEvent* event) {
    if (event->button() == Qt::LeftButton &&
        // Nicht dragging wenn im Overlay-Bereich
        !m_overlay->geometry().contains(event->pos()))
    {
        m_dragging  = true;
        m_dragStart = event->globalPosition().toPoint()
                      - frameGeometry().topLeft();
    }
    QMainWindow::mousePressEvent(event);
}

void MadoludusWindow::mouseMoveEvent(QMouseEvent* event) {
    if (m_dragging && (event->buttons() & Qt::LeftButton))
        move(event->globalPosition().toPoint() - m_dragStart);
    QMainWindow::mouseMoveEvent(event);
}
\end{lstlisting}

\verb|globalPosition()| liefert die absolute Bildschirmposition
der Maus (Qt6: gibt \verb|QPointF| zurueck, daher
\verb|.toPoint()|). Der Offset \member{m\_dragStart}
ist die Differenz zwischen Mausposition und Fenster-Ecke --
so bewegt sich das Fenster nicht sprunghaft.

\subsubsection{Qt-Analyse: FF-Flip-Book-Pipeline}

\begin{lstlisting}[caption={startFFExtraction: VideoFramesWorker starten}]
void MadoludusWindow::startFFExtraction(const QString& path)
{
    stopFFExtraction();

    const VideoFramesWorker::Mode mode =
        m_cfg.ffFixedCount
        ? VideoFramesWorker::Mode::FixedCount
        : VideoFramesWorker::Mode::IntervalSeconds;

    auto* worker = new VideoFramesWorker(
        path, mode, m_cfg.ffValue);
    m_ffTmpDir = worker->tmpDir();
    auto* thread = new QThread;
    m_ffThread = thread;
    worker->moveToThread(thread);

    connect(thread, &QThread::started,
            worker, &VideoFramesWorker::run);
    connect(worker, &VideoFramesWorker::resultReady,
            this,   &MadoludusWindow::onFFFramesReady);
    connect(worker, &VideoFramesWorker::resultReady,
            thread, &QThread::quit);
    connect(worker, &VideoFramesWorker::resultReady,
            worker, &QObject::deleteLater);
    connect(thread, &QThread::finished,
            thread, &QObject::deleteLater);
    thread->start();
}

void MadoludusWindow::onFFFramesReady(
    bool ok, QStringList framePaths)
{
    m_ffThread = nullptr;
    if (!ok || framePaths.isEmpty()) {
        m_display->setText(
            "(kein Video-Preview -- ffmpeg installieren)");
        startAutoAdvance();
        return;
    }
    m_ffFrames    = framePaths;
    m_ffFrameIdx  = 0;
    cycleFFFrame();
    if (!m_paused) {
        m_ffCycleTimer->start(kFFCycleMs);
        // Auto-Advance nach allen Frames einmal gesehen
        m_advTimer->start(
            qMax(kImageDelayFF,
                 kFFCycleMs * m_ffFrames.size()));
    }
}
\end{lstlisting}

\subsubsection{Qt-Analyse: scaleCurrentImage beim Resize}

\begin{lstlisting}[caption={resizeEvent: Bild neu skalieren}]
void MadoludusWindow::resizeEvent(QResizeEvent* event) {
    QMainWindow::resizeEvent(event);
    scaleCurrentImage();
}

void MadoludusWindow::scaleCurrentImage() {
    if (!m_display->isVisible() ||
        m_display->pixmap().isNull()) return;

    if (m_usingFlipbook && !m_ffFrames.isEmpty()) {
        // Letzten angezeigten Frame neu skalieren
        const int shown =
            (m_ffFrameIdx + m_ffFrames.size() - 1)
            % m_ffFrames.size();
        const QPixmap pm(m_ffFrames.at(shown));
        if (!pm.isNull())
            m_display->setPixmap(pm.scaled(
                m_display->size(),
                Qt::KeepAspectRatio,
                Qt::SmoothTransformation));
    } else if (!m_isVideo && m_index >= 0) {
        // Statisches Bild neu laden und skalieren
        QImageReader reader(m_mediaFiles.at(m_index));
        const QSize win = m_display->size();
        if (!win.isEmpty()) {
            reader.setScaledSize(win);
            const QImage img = reader.read();
            if (!img.isNull())
                m_display->setPixmap(
                    QPixmap::fromImage(img).scaled(
                        win, Qt::KeepAspectRatio,
                        Qt::SmoothTransformation));
        }
    }
}
\end{lstlisting}

% ----------------------------------------------------------
\subsection{madoludustab.h / madoludustab.cpp}
% ----------------------------------------------------------

\subsubsection{Klassendeklaration}

\begin{lstlisting}[style=header, caption={madoludustab.h -- Kern}]
class MadoludusTab : public QWidget {
    Q_OBJECT
public:
    explicit MadoludusTab(
        std::function<QString()> getGlobalDir,
        QWidget* parent = nullptr);
    ~MadoludusTab() override;  // stoppt ggf. Scan-Thread

    void onDirectoryChanged(const QString& path,
                            const QStringList& allFiles = {});
    void setSecretMode(bool on);

private slots:
    void showPreview(const QString& path);
    void onStartClicked();
    void onScanDone(bool ok,
                    QStringList imageFiles,
                    QStringList allFiles);
    void onWindowClosed();
    void openWindow(const QStringList& files);

private:
    void buildUi();
    void scanDir(const QString& path);
    MadoludusConfig buildConfig() const;
    QStringList filteredFiles() const;

    std::function<QString()> m_getGlobalDir;
    QStringList m_allFiles;
    bool        m_secretMode = false;

    QCheckBox*   m_chkImages  = nullptr;
    QCheckBox*   m_chkVideos  = nullptr;
    QCheckBox*   m_chkFF      = nullptr;
    QComboBox*   m_cmbFFRate  = nullptr;  // v0.5.42
    QCheckBox*   m_chkMuted   = nullptr;
    QCheckBox*   m_chkRandom  = nullptr;  // v0.5.42
    QPushButton* m_btnPreview = nullptr;
    QPushButton* m_btnStart   = nullptr;
    QLabel*      m_lblStatus  = nullptr;
    QProgressBar* m_pb        = nullptr;

    VideoPreviewWidget* m_preview = nullptr;
    QLabel*             m_lblType = nullptr;
    QThread*        m_scanThread  = nullptr;
    MadoludusWindow* m_window     = nullptr;
};
\end{lstlisting}

\subsubsection{C++-Analyse: buildConfig und FF-Rate-Tabelle}

\begin{lstlisting}[caption={FF-Rate-Tabelle und buildConfig}]
// Tabelle aller FF-Raten-Optionen:
struct FFRateEntry {
    const char* label;
    bool   fixedCount; // false=Intervall, true=Anzahl
    double value;
};
static const FFRateEntry kFFRates[] = {
    { "1/sec",       false, 1.0  },
    { "1/2 sec",     false, 2.0  },
    { "1/5 sec",     false, 5.0  },
    { "1/10 sec",    false, 10.0 },
    { "1/2 frames",  true,  2.0  },
    { "1/10 frames", true,  10.0 },
    { "1/30 frames", true,  30.0 },
    { "1/60 frames", true,  60.0 },
};
static constexpr int kFFRateCount =
    static_cast<int>(sizeof(kFFRates)/sizeof(kFFRates[0]));

// Config aus UI-Zustand bauen:
MadoludusConfig MadoludusTab::buildConfig() const {
    MadoludusConfig cfg;
    cfg.showImages  = m_chkImages->isChecked();
    cfg.showVideos  = m_chkVideos->isChecked();
    cfg.ffMode      = m_chkFF->isChecked();
    cfg.muted       = m_chkMuted->isChecked();
    cfg.randomMode  = m_chkRandom->isChecked();

    const int idx = m_cmbFFRate->currentIndex();
    if (idx >= 0 && idx < kFFRateCount) {
        cfg.ffFixedCount = kFFRates[idx].fixedCount;
        cfg.ffValue      = kFFRates[idx].value;
    } else {
        cfg.ffFixedCount = false;
        cfg.ffValue      = 10.0;
    }
    return cfg;
}
\end{lstlisting}

\verb|sizeof(kFFRates)/sizeof(kFFRates[0])| berechnet
die Arraylaenge zur Compile-Zeit ohne hartcodierten Wert --
bei jedem neuen Eintrag in der Tabelle stimmt der Count
automatisch.

\subsubsection{Qt-Analyse: openWindow und Signal-Rueckkanal}

\begin{lstlisting}[caption={openWindow: Fenster starten und verbinden}]
void MadoludusTab::openWindow(const QStringList& files) {
    MadoludusConfig cfg = buildConfig();
    cfg.secretMode = m_secretMode;
    m_window = new MadoludusWindow(files, cfg);

    // Startgroesse: halber Bildschirm, linke obere Ecke
    const QRect scr =
        QGuiApplication::primaryScreen()
        ->availableGeometry();
    m_window->setGeometry(
        0, 0, scr.width()/2, scr.height()/2);

    // Rueckkanal: Fenster meldet geschlossen
    connect(m_window, &MadoludusWindow::windowClosed,
            this,     &MadoludusTab::onWindowClosed);

    // Rueckkanal: Datei geaendert -> Preview aktualisieren
    connect(m_window,
            &MadoludusWindow::currentFileChanged,
            this,
            [this](const QString& path) {
                showPreview(path);
            });

    m_window->show();
    m_lblStatus->setText(
        QString("MadoLudus running -- %1 file(s).")
        .arg(files.size()));
}

void MadoludusTab::onWindowClosed() {
    m_window = nullptr;  // Zeiger freigeben (WA_DeleteOnClose)
    m_lblStatus->setText("MadoLudus closed.");
}
\end{lstlisting}

\verb|WA_DeleteOnClose| (gesetzt im \klasse{MadoludusWindow}-Konstruktor)
loescht das Objekt beim Schliessen automatisch.
Der Zeiger \member{m\_window} wird im \slot{onWindowClosed}-Slot
auf \verb|nullptr| gesetzt -- sonst haette man einen
Dangling Pointer.

% ----------------------------------------------------------
\subsection{madooverlay.h / madooverlay.cpp}
% ----------------------------------------------------------

\klasse{MadoOverlay} ist die halbtransparente Steuerleiste
im \klasse{MadoludusWindow}. Sie ist ein eigenstaendiges
\klasse{QWidget}, das als Kind im Fenster positioniert wird.

\begin{lstlisting}[style=header, caption={madooverlay.h}]
class MadoOverlay : public QWidget {
    Q_OBJECT
public:
    explicit MadoOverlay(QWidget* parent = nullptr);

    void setPaused(bool paused);   // Play/Pause-Icon
    void setMuted(bool muted);     // Lautsprecher-Icon
    void setFFMode(bool ff);       // FF-Button-Farbe
    void setRandomMode(bool on);   // v0.5.42: Shuffle-Icon
    void setFileInfo(int index, int total,
                     const QString& name);
    void setSecretMode(bool on);

signals:
    void playPauseClicked();
    void prevClicked();
    void nextClicked();
    void muteClicked();
    void ffClicked();
    void shuffleClicked();         // v0.5.42

protected:
    void paintEvent(QPaintEvent* event) override;

private:
    QPushButton* m_btnPrev    = nullptr;
    QPushButton* m_btnPlay    = nullptr;
    QPushButton* m_btnNext    = nullptr;
    QPushButton* m_btnMute    = nullptr;
    QPushButton* m_btnFF      = nullptr;
    QPushButton* m_btnShuffle = nullptr;  // v0.5.42
    QLabel*      m_lblInfo    = nullptr;
};
\end{lstlisting}

\subsubsection{Qt-Analyse: Custom paintEvent}

\begin{lstlisting}[caption={MadoOverlay::paintEvent -- dunkle abgerundete Pille}]
void MadoOverlay::paintEvent(QPaintEvent*) {
    QPainter p(this);
    p.setRenderHint(QPainter::Antialiasing);
    // Schwarz mit 63% Deckkraft (Alpha 160/255)
    p.setBrush(QColor(0, 0, 0, 160));
    p.setPen(Qt::NoPen);
    p.drawRoundedRect(
        rect().adjusted(4, 4, -4, -4), 10, 10);
}
\end{lstlisting}

\verb|QColor(r, g, b, alpha)| -- der vierte Parameter
ist der Alpha-Kanal (0 = voellig transparent,
255 = voellig opak). \verb|drawRoundedRect| mit Radius 10
erzeugt die "Pille"-Form. \verb|Qt::NoPen| verhindert
einen Rahmen.

\begin{hinweis}
\textbf{Wann paintEvent ueberschreiben?}
Normalerweise reicht es, ein \klasse{QWidget} per
\verb|setStyleSheet()| zu stylen.
\verb|paintEvent| ueberschreibt man nur, wenn man
volle Kontrolle ueber das Zeichnen braucht -- hier
fuer die halbtransparente Pille, die kein
Stylesheet erreichen kann.
\end{hinweis}

\subsubsection{Qt-Analyse: Unicode-Icons in Buttons}

\begin{lstlisting}[caption={makeBtn: Emoji-Symbole als Button-Texte}]
static QPushButton* makeBtn(const QString& text,
                             QWidget* parent) {
    auto* b = new QPushButton(text, parent);
    b->setFixedSize(40, 40);
    b->setFlat(true);
    b->setStyleSheet(
        "QPushButton {"
        "  font-size: 18px;"
        "  color: white;"
        "  background: transparent;"
        "  border: none;"
        "  border-radius: 6px;"
        "}"
        "QPushButton:hover {"
        "  background: rgba(255,255,255,40);"
        "}");
    return b;
}
// Aufruf:
m_btnPlay    = makeBtn("\u23f8", this);  // Pause-Symbol
m_btnMute    = makeBtn("\U0001f507", this); // Stummschalten
m_btnShuffle = makeBtn("\U0001f500", this); // Shuffle
\end{lstlisting}

Qt stellt Unicode-Zeichen als String-Literale dar.
\verb|\uXXXX| (4 Hex-Ziffern) fuer BMP-Zeichen,
\verb|\UXXXXXXXX| (8 Hex-Ziffern) fuer Zeichen
ausserhalb der BMP (Supplementary Planes, z.B. Emojis).
Die Schriftart muss die Zeichen unterstuetzen --
auf Windows liefert "Segoe UI Emoji" alle gaengigen Emojis.

% ============================================================
%  Kapitel 5: Worker-Klassen und gemeinsame Infrastruktur
% ============================================================
\chapter{Worker-Klassen und gemeinsame Infrastruktur}

% ============================================================
\section{Funktionsweise}
% ============================================================

TilVista fuehrt alle datei-intensiven Operationen in
Hintergrund-Threads aus, damit der GUI-Thread jederzeit
reaktiv bleibt. Alle Worker folgen demselben Muster:

\begin{enumerate}[noitemsep]
  \item \klasse{QObject}-Subklasse mit einem \verb|run()|-Slot
  \item \verb|moveToThread()| in einen \klasse{QThread}
  \item Arbeit in \methode{run}, Ergebnis per Signal
  \item Selbstaufraeuumung via \verb|deleteLater()|
\end{enumerate}

Die Worker-Klassen im Ueberblick:
\begin{center}
\small
\begin{tabular}{lll}
\toprule
\textbf{Klasse} & \textbf{Aufgabe} & \textbf{Genutzt von} \\
\midrule
\klasse{ScanWorker}       & Verzeichnis rekursiv scannen  &
  AleaVueTab, MadoludusTab \\
\klasse{DBSaveWorker}     & JSON-Datei schreiben          &
  DirDatabasePanel, ShujukoPanel \\
\klasse{DBLoadWorker}     & JSON-Datei lesen              &
  DirDatabasePanel, ShujukoPanel \\
\klasse{CatalogueWriteWorker} & .dshow/.catalogue schreiben &
  DirDatabasePanel \\
\klasse{CatalogueLoadWorker}  & .dshow/.catalogue lesen     &
  DirDatabasePanel \\
\klasse{VideoFramesWorker}    & ffmpeg-Frames extrahieren   &
  VideoPreviewWidget, MadoludusWindow \\
\bottomrule
\end{tabular}
\end{center}

% ============================================================
\section{File by File}
% ============================================================

\subsection{scanworker.h / scanworker.cpp}

\begin{lstlisting}[style=header, caption={scanworker.h}]
class ScanWorker : public QObject {
    Q_OBJECT
public:
    explicit ScanWorker(const QString& directory,
                        QObject* parent = nullptr);
public slots:
    void run();
signals:
    void progressChanged(int percent);
    void resultReady(bool ok,
                     QStringList imageFiles,
                     QStringList allFiles);
private:
    QString m_directory;
};
\end{lstlisting}

\subsubsection{C++-Analyse: run() -- rekursiver Verzeichnis-Scan}

\begin{lstlisting}[caption={ScanWorker::run -- QDirIterator mit Subdirectories}]
void ScanWorker::run() {
    QStringList all, img;
    try {
        // Rekursive Iteration ohne Limit
        QDirIterator it(
            m_directory,
            QDir::Files | QDir::NoDotAndDotDot,
            QDirIterator::Subdirectories);

        while (it.hasNext()) {
            it.next();
            all << it.filePath();
        }
        const int total = all.isEmpty() ? 1 : all.size();
        const auto& suf = TV::imageSuffixes();

        for (int i = 0; i < all.size(); ++i) {
            const QString ext =
                '.' + QFileInfo(all[i]).suffix().toLower();
            if (suf.contains(ext))
                img << all[i];
            // Fortschritt in den GUI-Thread senden
            emit progressChanged((i+1)*100/total);
        }
        emit resultReady(true, img, all);
    } catch (...) {
        // Jede Exception abfangen: im Worker-Thread gibt
        // es keinen Aufrufer, der Exceptions behandeln koennte
        emit resultReady(false, {}, {});
    }
}
\end{lstlisting}

\begin{qtinfo}[title={Qt-Hintergrund: QDirIterator}]
\klasse{QDirIterator} traversiert einen Verzeichnisbaum
iterativ ohne Rekursion -- kein Stack-Overflow-Risiko auch
bei tief verschachtelten Strukturen.
Die Option \verb|Subdirectories| aktiviert die rekursive
Traversal. \verb|NoDotAndDotDot| schliesst \datei{.} und
\datei{..} aus.\\
Referenz: \url{https://doc.qt.io/qt-6/qdiriterator.html}
\end{qtinfo}

\paragraph{try/catch(...) im Worker-Thread}
Der \verb|catch(...)| (drei Punkte = jede Exception)
ist hier bewusst breit: Im Worker-Thread gibt es keinen
uebergeordneten Aufrufer auf dem Call-Stack, der eine
Exception sinnvoll behandeln koennte. Unbehandelte
Exceptions in einem Thread beenden das Programm unkontrolliert
(\verb|std::terminate|). Daher: abfangen, \verb|false|
melden, sauber beenden.

% ----------------------------------------------------------
\subsection{dbworkers.h / dbworkers.cpp}
% ----------------------------------------------------------

\begin{lstlisting}[style=header, caption={dbworkers.h}]
class DBSaveWorker : public QObject {
    Q_OBJECT
public:
    DBSaveWorker(const QString& path,
                 const QJsonObject& data,
                 QObject* p = nullptr);
public slots: void run();
signals: void resultReady(bool ok);
private:
    QString     m_path;
    QJsonObject m_data;  // Wert-Kopie (thread-sicher)
};

class DBLoadWorker : public QObject {
    Q_OBJECT
public:
    explicit DBLoadWorker(const QString& path,
                          QObject* p = nullptr);
public slots: void run();
signals: void resultReady(bool ok, QJsonObject data);
private:
    QString m_path;
};
\end{lstlisting}

\subsubsection{C++-Analyse: QJsonObject als Wert-Kopie}

\begin{lstlisting}[caption={DBSaveWorker::run}]
void DBSaveWorker::run() {
    QDir().mkpath(QFileInfo(m_path).absolutePath());
    QFile f(m_path);
    if (!f.open(QIODevice::WriteOnly | QIODevice::Truncate)) {
        emit resultReady(false);
        return;
    }
    // QJsonDocument::Indented: lesbare Formatierung
    f.write(QJsonDocument(m_data)
            .toJson(QJsonDocument::Indented));
    emit resultReady(true);
}
\end{lstlisting}

\member{m\_data} ist eine \emph{Wert-Kopie} von
\verb|QJsonObject|. Qt-Wertobjekte nutzen
\emph{Copy-on-Write (Implicit Sharing)}: Der Kopiervorgang
im Konstruktor ist sehr guenstig (nur Referenzzaehler
erhoehen). Erst wenn jemand schreibt, wird wirklich kopiert.
Ergebnis: Beide Threads (GUI und Worker) haben
unabhaengige Daten ohne Mutex.

\verb|QDir().mkpath(...)| erzeugt das Zielverzeichnis
rekursiv, falls es noch nicht existiert --
das Aequivalent zu \texttt{mkdir -p}.

% ----------------------------------------------------------
\subsection{catalogueworkers.h / catalogueworkers.cpp}
% ----------------------------------------------------------

\begin{lstlisting}[style=header, caption={catalogueworkers.h}]
class CatalogueWriteWorker : public QObject {
    Q_OBJECT
public:
    CatalogueWriteWorker(
        const QString& kaivoDir,
        const QString& dbPath,
        const QJsonObject& dbData,
        const QString& entryName,
        const QStringList& imageFiles,
        const QStringList& allFiles,
        QObject* p = nullptr);
public slots: void run();
signals: void resultReady(bool ok);
};

class CatalogueLoadWorker : public QObject {
    Q_OBJECT
public:
    CatalogueLoadWorker(
        const QString& kaivoDir,
        const QString& imgFname,
        const QString& allFname,
        QObject* p = nullptr);
public slots: void run();
signals:
    void resultReady(bool ok,
                     QStringList imageFiles,
                     QStringList allFiles);
private:
    QString m_kaivoDir, m_imgFname, m_allFname;
    // static: kein this-Pointer noetig
    static QStringList readLines(const QString& path);
};
\end{lstlisting}

\subsubsection{C++-Analyse: static readLines}

\verb|readLines| ist eine statische Hilfsmethode:
Sie braucht keinen Zustand des Objekts, nur einen Pfad.
\verb|static| im Klassenkontext bedeutet: kein impliziter
\verb|this|-Zeiger, kein Zugriff auf Instanz-Member.
Man haette auch eine freie Funktion nehmen koennen --
als statische Methode ist sie jedoch logisch der Klasse
zugeordnet.

\begin{lstlisting}[caption={CatalogueLoadWorker::readLines}]
QStringList CatalogueLoadWorker::readLines(
    const QString& path)
{
    QStringList r;
    QFile f(path);
    if (!f.open(QIODevice::ReadOnly | QIODevice::Text))
        return r;
    QTextStream in(&f);
    while (!in.atEnd()) {
        const QString l = in.readLine().trimmed();
        if (!l.isEmpty()) r << l;
    }
    return r;
}
\end{lstlisting}

\verb|QTextStream| liest zeilenweise und behandelt
Encoding-Fragen. \verb|trimmed()| entfernt fuehrende
und abschliessende Leerzeichen (inkl. \verb|\r| unter
Windows) -- wichtig fuer Cross-Platform-Kompatibilitaet
der Katalog-Dateien.

% ----------------------------------------------------------
\subsection{videoframesworker.h / videoframesworker.cpp}
% ----------------------------------------------------------

\klasse{VideoFramesWorker} extrahiert Einzelbilder aus
Videos via ffmpeg und legt sie als JPEG im Temp-Verzeichnis
ab.

\begin{lstlisting}[style=header, caption={videoframesworker.h -- Kern}]
class VideoFramesWorker : public QObject {
    Q_OBJECT
public:
    enum class Mode { IntervalSeconds, FixedCount };

    explicit VideoFramesWorker(
        const QString& videoPath,
        Mode    mode  = Mode::IntervalSeconds,
        double  value = 10.0,
        QObject* parent = nullptr);

    const QString& tmpDir() const { return m_tmpDir; }

public slots: void run();

signals:
    void resultReady(bool ok, QStringList framePaths);

private:
    // Videodauer via "ffmpeg -i" aus stderr lesen
    static double probeDurationSeconds(const QString& path);
    QList<double> computeTimestamps(double durationSeconds) const;

    QString m_videoPath, m_tmpDir;
    Mode    m_mode;
    double  m_value;
};
\end{lstlisting}

\subsubsection{C++-Analyse: Scoped Enum (enum class)}

\begin{lstlisting}
enum class Mode { IntervalSeconds, FixedCount };
\end{lstlisting}

\verb|enum class| (C++11) ist ein stark typisiertes Enum:
\begin{itemize}[noitemsep]
  \item Werte sind \emph{nicht} in den umgebenden Scope
    exportiert: Man schreibt
    \verb|Mode::IntervalSeconds|, nicht einfach
    \verb|IntervalSeconds|
  \item Kein implizites Konvertieren zu \verb|int|:
    Fehler werden vom Compiler erkannt statt still
    falsch zu laufen
  \item Besser lesbar als \verb|bool ffFixedCount|
    mit zwei Bedeutungen
\end{itemize}

\subsubsection{C++-Analyse: probeDurationSeconds}

\begin{lstlisting}[caption={probeDurationSeconds: ffmpeg-Stderr parsen}]
double VideoFramesWorker::probeDurationSeconds(
    const QString& path)
{
    // "ffmpeg -i <path>" gibt Metadaten auf stderr aus
    // und beendet sich mit Fehlercode (kein Output-File).
    // Genau das brauchen wir -- nur die Metadaten.
    QProcess proc;
    proc.start("ffmpeg", {"-i", path});
    proc.waitForFinished(8000);
    const QString out =
        QString::fromUtf8(proc.readAllStandardError());

    // "Duration: HH:MM:SS.xx" im stderr suchen
    static const QRegularExpression re(
        R"(Duration:\s*(\d+):(\d+):(\d+(?:\.\d+)?))");
    const auto m = re.match(out);
    if (!m.hasMatch()) return 0.0;

    return m.captured(1).toDouble() * 3600.0  // Stunden
         + m.captured(2).toDouble() *   60.0  // Minuten
         + m.captured(3).toDouble();           // Sekunden
}
\end{lstlisting}

\textbf{Warum stderr?} ffmpeg gibt Metadaten immer auf
\emph{stderr} aus, auch wenn es sich um normale
Informationen handelt. stdout ist fuer die eigentliche
Ausgabedatei reserviert.

\textbf{Raw-String-Regex}: \verb|R"(...)"| vermeidet
doppeltes Escaping. Ohne Raw-String:
\verb|"Duration:\\s*(\\d+):(\\d+):"| -- schlecht lesbar.

\subsubsection{v0.5.42: computeTimestamps -- durationssensitiv}

\begin{lstlisting}[caption={computeTimestamps: IntervalSeconds und FixedCount}]
QList<double> VideoFramesWorker::computeTimestamps(
    double durationSeconds) const
{
    constexpr int kMaxFrames = 60;
    QList<double> ts;

    if (durationSeconds <= 0.5) {
        ts << 0.0;
        return ts;
    }

    if (m_mode == Mode::IntervalSeconds) {
        // Ein Frame pro "value" Sekunden
        const double step = qMax(0.2, m_value);
        for (double t = 0.0; t < durationSeconds; t += step)
            ts << t;
        if (ts.isEmpty()) ts << 0.0;
    } else {
        // Genau "value" Frames gleichmaessig verteilt
        const int count = qMax(1, int(m_value));
        if (count == 1) {
            ts << durationSeconds / 2.0;
            return ts;
        }
        for (int i = 0; i < count; ++i)
            ts << durationSeconds * double(i)
                                  / double(count - 1);
    }

    // Auf 60 Frames deckeln (verhindert hunderte ffmpeg-Aufrufe)
    if (ts.size() > kMaxFrames) {
        QList<double> capped;
        for (int i = 0; i < kMaxFrames; ++i)
            capped << ts.at(
                i * (ts.size()-1) / (kMaxFrames-1));
        return capped;
    }
    return ts;
}
\end{lstlisting}

\begin{neuin}[title={Neu in v0.5.42: Dauer-bewusste Timestamps}]
Fruehre Version: fest \{0, 10, 20, 30, 40, 50\} Sekunden --
bei kurzen Videos (unter 10 Sekunden) gab es fast keine
Frames, bei langen Videos immer nur den Anfang.

Neue Version: Die Videodauer wird per \verb|ffmpeg -i|
ermittelt. Dann werden Timestamps relativ zur Gesamtdauer
berechnet -- fuer ein 2-Minuten-Video mit "1/sec" gibt
es 120 Frames (gedeckelt auf 60). Fuer "1/10 frames" gibt
es immer genau 10 gleichmaessig verteilte Frames.
\end{neuin}

\subsubsection{Qt-Analyse: QProcess fuer externe Programme}

\begin{lstlisting}[caption={run(): ffmpeg-Frames extrahieren}]
void VideoFramesWorker::run() {
    // ffmpeg vorhanden?
    if (QStandardPaths::findExecutable("ffmpeg").isEmpty()) {
        emit resultReady(false, {});
        return;
    }
    const double duration = probeDurationSeconds(m_videoPath);
    const QList<double> times = duration > 0.0
        ? computeTimestamps(duration)
        : QList<double>{0.0, 5.0, 10.0, 15.0, 20.0, 25.0};

    QStringList frames;
    int idx = 0;
    for (double t : times) {
        const QString out =
            QString("%1/f%2.jpg")
            .arg(m_tmpDir)
            // 3-stellig, Basis 10, Fuellzeichen '0'
            .arg(idx++, 3, 10, QChar('0'));

        QProcess proc;
        proc.start("ffmpeg", {
            "-i",        m_videoPath,
            "-ss",       QString::number(t, 'f', 2),
            "-vframes",  "1",
            "-f",        "image2",
            "-y",        out
        });
        proc.waitForFinished(12000);
        if (proc.exitCode() == 0 && QFile::exists(out))
            frames << out;
    }
    emit resultReady(!frames.isEmpty(), frames);
}
\end{lstlisting}

\paragraph{QString::arg mit Formatierung}
\verb|.arg(idx++, 3, 10, QChar('0'))| formatiert
eine Zahl mit:
\begin{itemize}[noitemsep]
  \item Breite 3 (\verb|3|)
  \item Basis 10 (\verb|10|)
  \item Fuellzeichen \verb|'0'|
\end{itemize}
Ergebnis: 0 $\to$ \verb|f000.jpg|, 10 $\to$ \verb|f010.jpg|.
Sortierbare Dateinamen ohne Extrabibliothek.

\begin{qtinfo}[title={Qt-Hintergrund: QProcess und Argumente}]
\klasse{QProcess} startet externe Programme sicher ohne
Shell-Interaktion. \verb|proc.start(programm, arglist)|
uebergibt Argumente als \klasse{QStringList} -- Leerzeichen
in Pfaden werden korrekt behandelt, Shell-Injection ist
nicht moeglich. Im Gegensatz zu
\verb|system("ffmpeg -i " + path)| ist das auch auf
Pfaden mit Leerzeichen und Sonderzeichen zuverlaessig.\\
Referenz: \url{https://doc.qt.io/qt-6/qprocess.html}
\end{qtinfo}

% ----------------------------------------------------------
\subsection{Tab 4 -- shortcutstab.h / shortcutstab.cpp}
% ----------------------------------------------------------

Der About-Tab ist ein reiner Informations-Tab ohne
Datenbankanbindung.

\subsubsection{Datentabellen als statische Structs}

\begin{lstlisting}[caption={Datenmodell als Compile-Zeit-Konstanten}]
struct Row     { const char* key; const char* description; };
struct Section { const char* title; QList<Row> rows; };

static const QList<Section> kSections = {
    {
        "Global",
        {
            {"Ctrl+Alt+F8", "Toggle Secret Mode"},
            {"Enter (DirBar)", "Verzeichnis uebernehmen"},
        }
    },
    {
        "AleaVue  --  Vollbild-Diashow",
        {
            {"->",      "Naechstes Bild (zufaellig)"},
            {"<-",      "Zurueck (Verlauf)"},
            {"S",       "Bookmark in shujuko"},
            {"Space",   "Timer pause/resume"},
            {"Esc",     "Diashow schliessen"},
        }
    },
    // ...
};
\end{lstlisting}

\verb|const char*| ist ein C-String-Zeiger direkt auf
das String-Literal im Programm-Daten-Segment. Da die
Strings nicht geaendert werden, ist das effizient und
sicher. \klasse{QList<Row>} kann mit
Initializer-Listen \verb|{ {...}, {...} }| befuellt werden.

\subsubsection{Dynamische Tabellenhoehe}

\begin{lstlisting}[caption={makeTable: exakte Hoehe berechnen}]
static QTableWidget* makeTable(const Section& sec) {
    auto* table = new QTableWidget(
        static_cast<int>(sec.rows.size()), 2);
    table->horizontalHeader()
         ->setStretchLastSection(true);
    table->verticalHeader()->setVisible(false);
    table->setEditTriggers(
        QAbstractItemView::NoEditTriggers);
    table->setSelectionMode(
        QAbstractItemView::NoSelection);
    table->setShowGrid(false);
    table->setAlternatingRowColors(true);

    // ... Eintraege setzen ...
    table->resizeRowsToContents();

    // Exakte Hoehe = Headerhoehe + alle Zeilenhoehen
    int h = table->horizontalHeader()->height();
    for (int r = 0; r < table->rowCount(); ++r)
        h += table->rowHeight(r);
    table->setFixedHeight(h + 4);  // +4 Pixel Puffer
    return table;
}
\end{lstlisting}

\verb|setFixedHeight| klemmt die Hoehe auf exakt den
Inhalt -- kein Leerraum, keine Scrollbar noetig.
\verb|setAlternatingRowColors(true)| erzeugt abwechselnde
Zeilenfaerbung automatisch -- per OS-Theme-Farben,
kein Hardcoding noetig.

% ============================================================
%  Kapitel 6: Zusammenfassung
% ============================================================
\chapter{Zusammenfassung}

% ============================================================
\section{Architektur -- Rueckblick}
% ============================================================

TilVista v0.5.42 demonstriert eine durchgaengige,
signalgesteuerte Qt6-Architektur. Die wichtigsten
Entwurfsentscheidungen im Rueckblick:

\subsection{Entkopplung durch std::function-Callbacks}

Kein Tab importiert den Header eines anderen Tabs.
Statt \verb|MainWindow*| als Parameter erhalten alle
Tabs ein Lambda:

\begin{lstlisting}
// In MainWindow::MainWindow():
m_aleaVueTab = new AleaVueTab(
    [this]{ return m_dirBar->directory(); },
    m_shujuko);
\end{lstlisting}

Das entspricht dem \emph{Dependency Inversion Principle}
(DIP) aus den SOLID-Prinzipien: Abhaengigkeit von einer
Abstraktion (Funktion), nicht von einer konkreten Klasse.

\subsection{Unidirektionale Signalkette}

Zustandsaenderungen fliessen immer von einem Sender
durch \klasse{MainWindow} zu den betroffenen Empfaengern
-- niemals direkt zwischen Tabs:

\begin{center}
\small
\begin{tabular}{ccc}
\texttt{DirDatabasePanel} &
\(\longrightarrow\) \textit{Signal} \(\longrightarrow\) &
\texttt{MainWindow} \\
& & \(\downarrow\) \\
& & \texttt{SattumaPicTab} \\
& & \texttt{MadoludusTab} \\
\end{tabular}
\end{center}

Kein Tab reagiert direkt auf Signale eines anderen Tabs.
Das macht den Datenfluss nachvollziehbar und
Fehlerquellen leicht lokalisierbar.

\subsection{Worker-Thread-Muster durchgaengig}

Jede datei-IO-Operation laeuft in einem eigenen
\klasse{QThread} mit dem Worker-Object-Muster:

\begin{lstlisting}
auto* worker = new MeinWorker(parameter);
auto* thread = new QThread;
worker->moveToThread(thread);
connect(thread, &QThread::started,
        worker, &MeinWorker::run);
connect(worker, &MeinWorker::resultReady,
        this,   &MeineKlasse::onErgebnis);
connect(worker, &MeinWorker::resultReady,
        thread, &QThread::quit);
connect(worker, &MeinWorker::resultReady,
        worker, &QObject::deleteLater);
connect(thread, &QThread::finished,
        thread, &QObject::deleteLater);
thread->start();
\end{lstlisting}

Der GUI-Thread blockiert nie. Fortschritt kommt per
\signal{progressChanged(int)}, Ergebnis per
\signal{resultReady}.

\subsection{safeStop-Idiom}

Vor jedem Thread-Start werden alle Signal-Verbindungen
des alten Threads zum Empfaenger getrennt:

\begin{lstlisting}
static void safeStop(QThread*& t, QObject* receiver) {
    if (!t) return;
    t->disconnect(receiver);
    t = nullptr;
}
\end{lstlisting}

Das verhindert Stale-Signal-Bugs ohne Mutex oder
\verb|QThread::wait()| (das wegen Queued-Connections
deadlocken wuerde).

\subsection{Optionale Abhaengigkeit: TV\_NO\_MULTIMEDIA}

Qt Multimedia wird per \cmake{find\_package QUIET}
gesucht. Nicht gefunden: \makro{TV\_NO\_MULTIMEDIA}
wird als Praeprozessor-Makro definiert. Im Code
steuern \verb|#ifndef|/\verb|#ifdef| die Codepfade.
Das Programm kompiliert und laeuft in beiden Faellen.

\subsection{v0.5.42: MadoLudus-Refactor}

Die Trennung von \klasse{MadoludusTab}
(Konfiguration) und \klasse{MadoludusWindow}
(Wiedergabe) reduzierte eine 800-Zeilen-Monolith-Klasse
auf zwei klar abgegrenzte Verantwortlichkeiten.
Das Fenster kommuniziert per zwei Signalen zurueck:
\signal{windowClosed()} und
\signal{currentFileChanged(QString)}.

% ============================================================
\section{C++-Konzepte im Ueberblick}
% ============================================================

\begin{longtable}{p{4.2cm}p{4.2cm}p{5cm}}
\toprule
\textbf{Konzept} & \textbf{Datei(en)} & \textbf{Zweck} \\
\midrule
\endfirsthead
\midrule
\endhead
\verb|std::function<T>| &
  aleavuetab, sattumapictab, madoludustab &
  Callback ohne harte Abhaengigkeit \\[4pt]

\verb|std::mt19937| + \verb|std::random_device| &
  sattumapictab, slideshowwindow, madoluduswindow &
  Qualitaetszufall, OS-Entropie als Seed \\[4pt]

\verb|static| (lokal) &
  pathutils, DirDatabasePanel &
  Einmalige Initialisierung, Singleton-Muster \\[4pt]

\verb|static constexpr| &
  slideshowwindow, videopreviewwidget, madoluduswindow &
  Compile-Zeit-Konstanten, typsicher \\[4pt]

\verb|= nullptr| / \verb|= false| (In-class-Init) &
  alle Klassen &
  Sichere Member-Initialisierung ohne Konstruktor \\[4pt]

\verb|override| &
  alle Qt-Subklassen &
  Compiler prueft korrekte Methodensignatur \\[4pt]

Raw-String \verb|R"(...)"| &
  pathutils, videoframesworker &
  Sonderzeichen ohne Escape-Sequenzen \\[4pt]

\verb|enum class| &
  videoframesworker &
  Stark typisiertes Enum, kein implizites \verb|int| \\[4pt]

\verb|try/catch(...)| &
  scanworker &
  Alle Exceptions im Worker-Thread abfangen \\[4pt]

Template-Funktion &
  dirdatabasepanel &
  DRY: generischer Thread-Starter \verb|launchWorker<>| \\[4pt]

\verb|sizeof(arr)/sizeof(arr[0])| &
  madoludustab &
  Compile-Zeit-Arraylaenge \\[4pt]

\verb|explicit| Konstruktor &
  alle Klassen &
  Verhindert ungewollte implizite Konvertierungen \\[4pt]

\verb|QStringLiteral| &
  pathutils &
  Compile-Zeit-String, schneller als Laufzeit-Konstruktion \\[4pt]

\verb|struct| mit Defaults &
  madoluduswindow &
  Konfigurationsdaten ohne Konstruktor-Boilerplate \\[4pt]

\bottomrule
\end{longtable}

% ============================================================
\section{Qt-Konzepte im Ueberblick}
% ============================================================

\begin{longtable}{p{4.0cm}p{4.4cm}p{5cm}}
\toprule
\textbf{Klasse / Konzept} & \textbf{Datei(en)} & \textbf{Zweck} \\
\midrule
\endhead

\klasse{QMainWindow} &
  mainwindow, slideshowwindow, madoluduswindow &
  Vollstaendiges Fenster mit Status\-leiste \\[4pt]

\klasse{QWidget} &
  alle Tab-Klassen, Panels &
  Basis aller GUI-Elemente \\[4pt]

\klasse{QTabWidget} &
  mainwindow &
  Tab-Navigation zwischen den Werkzeugen \\[4pt]

\klasse{QStackedWidget} &
  videopreviewwidget &
  Zwischen Seiten umschalten (Video / Bild) \\[4pt]

\klasse{QTimer} &
  slideshowwindow, madoluduswindow, videopreview &
  Zeitgesteuerte Ereignisse ohne Blocker \\[4pt]

\klasse{QThread} + moveToThread &
  alle Worker &
  Hintergrundarbeit ohne GUI-Blockade \\[4pt]

Queued Connection &
  alle Worker-Signale &
  Automatisch thread-sicherer Signal-Transport \\[4pt]

\verb|deleteLater()| &
  alle Worker &
  Thread-sicheres Aufraeuimen von QObjects \\[4pt]

\klasse{QDirIterator} &
  scanworker &
  Rekursive Verzeichnis-Traversal \\[4pt]

\klasse{QImageReader} &
  slideshowwindow, sattumapictab, madoluduswindow &
  Effizientes Laden mit Vorab-Skalierung \\[4pt]

\klasse{QPainter} &
  madooverlay, previewlabel &
  2D-Zeichenoperationen, Custom-Painting \\[4pt]

\klasse{QJsonObject} / \klasse{QJsonDocument} &
  config, dbworkers, shujukopanel &
  JSON-Persistenz, Copy-on-Write \\[4pt]

\klasse{QProcess} &
  videoframesworker &
  Externe Programme (ffmpeg) starten \\[4pt]

\klasse{QRegularExpression} &
  videoframesworker, dirdatabasepanel &
  PCRE2-Regex, statisch kompiliert \\[4pt]

\klasse{QShortcut} &
  mainwindow &
  App-weite Tastenkuerzel \\[4pt]

\klasse{QFileIconProvider} &
  sattumapictab &
  Native OS-Datei-Icons \\[4pt]

\klasse{QMediaPlayer} + \klasse{QVideoWidget} &
  videopreviewwidget, madoluduswindow &
  Native Videowiedergabe \\[4pt]

\makro{Q\_OS\_WIN} / \makro{TV\_NO\_MULTIMEDIA} &
  pathutils, alle Multimedia-Klassen &
  Plattform- und Feature-bedingte Compilation \\[4pt]

\verb|Qt::ApplicationShortcut| &
  mainwindow &
  Kuerzel auch bei Fokus in Eingabefeld aktiv \\[4pt]

\verb|Qt::FramelessWindowHint| &
  madoluduswindow &
  Rahmenlos + manuell verschiebbar \\[4pt]

\verb|WA_DeleteOnClose| &
  madoluduswindow, slideshowwindow &
  Automatisches Loeschen beim Schliessen \\[4pt]

\bottomrule
\end{longtable}

% ============================================================
\section{Empfohlener Lernpfad}
% ============================================================

Fuer den Einstieg in TilVista als Lernprojekt empfiehlt
sich folgende Lesereihenfolge -- von einfach nach komplex:

\begin{enumerate}
  \item \textbf{main.cpp} --
    Einstiegspunkt, \klasse{QApplication}, \klasse{Config},
    erstes Fenster; zwei Dutzend Zeilen die das grosse Bild
    zeigen

  \item \textbf{pathutils.h + pathutils.cpp} --
    Namespace, statische lokale Variablen,
    \verb|#ifdef Q_OS_WIN|, Raw-Strings; keine Qt-Abhaengigkeit
    ausser \klasse{QString}

  \item \textbf{config.h + config.cpp} --
    Statische Membervariablen, \klasse{QJsonDocument},
    einfachstes moegliches Singleton

  \item \textbf{scanworker.h + scanworker.cpp} --
    Erstes vollstaendiges Worker-Thread-Beispiel,
    \klasse{QDirIterator}, Signal-Fortschritt,
    \verb|try/catch|

  \item \textbf{mainwindow.h + mainwindow.cpp} --
    \klasse{QMainWindow}, Layout, Lambda-Callbacks,
    Signalkette, \verb|closeEvent|

  \item \textbf{slideshowwindow.h + slideshowwindow.cpp} --
    \klasse{QTimer}, \verb|keyPressEvent|,
    \klasse{QImageReader} mit Vorab-Skalierung,
    Verlaufslogik, Vollbild-Cursor

  \item \textbf{dirdatabasepanel.h + dirdatabasepanel.cpp} --
    JSON-Persistenz, drei parallele Threads,
    \verb|safeStop|, \verb|setBusy|,
    \verb|flushPendingWrites|, Template-Helfer

  \item \textbf{shujukopanel.h + shujukopanel.cpp} --
    Zweites JSON-Schema, Bool-Guard,
    \verb|validateFiles|, \verb|flushPendingWrites|

  \item \textbf{videoframesworker.h + videoframesworker.cpp} --
    \klasse{QProcess}, \verb|enum class|,
    dauer-bewusste Timestamps, Regex auf stderr

  \item \textbf{madoluduswindow.h + madoluduswindow.cpp} --
    Rahmenlos, Drag-to-Move, History-Navigation,
    optionale Multimedia, \verb|WA\_DeleteOnClose|

  \item \textbf{madoludustab.h + madoludustab.cpp} --
    Konfiguration trennen von Ausfuehrung,
    \verb|buildConfig()|, \verb|sizeof|-Trick

  \item \textbf{madooverlay.h + madooverlay.cpp} --
    Custom \verb|paintEvent|, Unicode-Emojis,
    Stylesheet-Zustandsaenderung per Methode
\end{enumerate}

% ============================================================
\section{Weiterfuehrende Ressourcen}
% ============================================================

\begin{description}[leftmargin=3cm, style=nextline]
  \item[Qt-Dokumentation]
    \url{https://doc.qt.io/qt-6/}\\
    Vollstaendige API-Referenz, Beispiele und Guides

  \item[Signals \& Slots]
    \url{https://doc.qt.io/qt-6/signalsandslots.html}\\
    Ausfuehrliche Erklaerung beider Syntaxen

  \item[Threading-Guide]
    \url{https://doc.qt.io/qt-6/thread-basics.html}\\
    Worker-Object-Muster, QThread-Lebenszyklus

  \item[Model/View]
    \url{https://doc.qt.io/qt-6/model-view-programming.html}\\
    Naechster Schritt nach dem Widget-basierten Ansatz

  \item[cppreference C++17]
    \url{https://en.cppreference.com/}\\
    Lambda, std::function, constexpr, enum class

  \item[C++ Core Guidelines]
    \url{https://isocpp.github.io/CppCoreGuidelines/}\\
    Empfohlene Idiome (RAII, Rule of Five, etc.)

  \item[CMake-Dokumentation]
    \url{https://cmake.org/cmake/help/latest/}\\
    find\_package, Generator-Expressions, Funktionen

  \item[CTAN: babel-german]
    \url{https://ctan.org/pkg/babel-german}\\
    Installation des fehlenden deutschen Babel-Pakets
    (dann kann in \datei{packages.tex} \verb|\usepackage[ngerman]{babel}|
    aktiviert werden)
\end{description}


% ============================================================
%  Anhang
% ============================================================
\appendix

\chapter{Tastaturkuerzel-Referenz}

\begin{longtable}{p{4.5cm}p{3cm}p{6cm}}
\toprule
\textbf{Taste} & \textbf{Kontext} & \textbf{Aktion} \\
\midrule
\endfirsthead
\midrule
\endfirsthead
\midrule
\endhead
Ctrl+Alt+F8    & Global        & Secret Mode ein/aus \\
Enter          & DirBar        & Verzeichnis uebernehmen \\
\midrule
Pfeil rechts   & AleaVue       & Naechstes Bild (zufaellig) \\
Pfeil links    & AleaVue       & Zurueck (Verlauf) \\
Leertaste      & AleaVue       & Timer pause/resume \\
S              & AleaVue       & Bookmark in shujuko \\
Esc            & AleaVue       & Diashow schliessen \\
\midrule
Space          & MadoLudus     & Play / Pause \\
Pfeil rechts   & MadoLudus     & Naechste Datei \\
Pfeil links    & MadoLudus     & Vorherige Datei (Verlauf) \\
Shift+Rechts   & MadoLudus (Video) & +5 Sekunden (native) \\
Shift+Links    & MadoLudus (Video) & -5 Sekunden (native) \\
M              & MadoLudus     & Stummschalten \\
F              & MadoLudus     & FF-Modus (Secret Mode zum Deaktiv.) \\
R              & MadoLudus     & Zufalls-/Sequenz-Reihenfolge \\
Esc            & MadoLudus     & Fenster schliessen \\
\midrule
{--}           & SattumaPic    & Nur Maus / Buttons \\
\bottomrule
\end{longtable}

\chapter{Dateistruktur zur Laufzeit}

\begin{lstlisting}[language={}, numbers=none,
  caption={Laufzeitverzeichnis neben der .exe}]
TilVista.exe
optegnelser/
  kaivo/
    tilvista_dirs.json          <- kaivo-Datenbank (DB1)
    <entryname>_img.dshow       <- Bildpfade (zeilenweise)
    <entryname>_all.catalogue   <- alle Pfade (zeilenweise)
    <entryname>_shujuko.json    <- Lesezeichen (DB2)
  loadingerrors.log             <- optional (AleaVue)
tilvista_config.json            <- Einstellungen
\end{lstlisting}

% -- Index ---------------------------------------------------
\printindex

\end{document}
